\documentclass[11pt]{article}

    \usepackage[breakable]{tcolorbox}
    \usepackage{parskip} % Stop auto-indenting (to mimic markdown behaviour)
    

    % Basic figure setup, for now with no caption control since it's done
    % automatically by Pandoc (which extracts ![](path) syntax from Markdown).
    \usepackage{graphicx}
    % Keep aspect ratio if custom image width or height is specified
    \setkeys{Gin}{keepaspectratio}
    % Maintain compatibility with old templates. Remove in nbconvert 6.0
    \let\Oldincludegraphics\includegraphics
    % Ensure that by default, figures have no caption (until we provide a
    % proper Figure object with a Caption API and a way to capture that
    % in the conversion process - todo).
    \usepackage{caption}
    \DeclareCaptionFormat{nocaption}{}
    \captionsetup{format=nocaption,aboveskip=0pt,belowskip=0pt}

    \usepackage{float}
    \floatplacement{figure}{H} % forces figures to be placed at the correct location
    \usepackage{xcolor} % Allow colors to be defined
    \usepackage{enumerate} % Needed for markdown enumerations to work
    \usepackage{geometry} % Used to adjust the document margins
    \usepackage{amsmath} % Equations
    \usepackage{amssymb} % Equations
    \usepackage{textcomp} % defines textquotesingle
    % Hack from http://tex.stackexchange.com/a/47451/13684:
    \AtBeginDocument{%
        \def\PYZsq{\textquotesingle}% Upright quotes in Pygmentized code
    }
    \usepackage{upquote} % Upright quotes for verbatim code
    \usepackage{eurosym} % defines \euro

    \usepackage{iftex}
    \ifPDFTeX
        \usepackage[T1]{fontenc}
        \IfFileExists{alphabeta.sty}{
              \usepackage{alphabeta}
          }{
              \usepackage[mathletters]{ucs}
              \usepackage[utf8x]{inputenc}
          }
    \else
        \usepackage{fontspec}
        \usepackage{unicode-math}
    \fi

    \usepackage{fancyvrb} % verbatim replacement that allows latex
    \usepackage{grffile} % extends the file name processing of package graphics
                         % to support a larger range
    \makeatletter % fix for old versions of grffile with XeLaTeX
    \@ifpackagelater{grffile}{2019/11/01}
    {
      % Do nothing on new versions
    }
    {
      \def\Gread@@xetex#1{%
        \IfFileExists{"\Gin@base".bb}%
        {\Gread@eps{\Gin@base.bb}}%
        {\Gread@@xetex@aux#1}%
      }
    }
    \makeatother
    \usepackage[Export]{adjustbox} % Used to constrain images to a maximum size
    \adjustboxset{max size={0.9\linewidth}{0.9\paperheight}}

    % The hyperref package gives us a pdf with properly built
    % internal navigation ('pdf bookmarks' for the table of contents,
    % internal cross-reference links, web links for URLs, etc.)
    \usepackage{hyperref}
    % The default LaTeX title has an obnoxious amount of whitespace. By default,
    % titling removes some of it. It also provides customization options.
    \usepackage{titling}
    \usepackage{longtable} % longtable support required by pandoc >1.10
    \usepackage{booktabs}  % table support for pandoc > 1.12.2
    \usepackage{array}     % table support for pandoc >= 2.11.3
    \usepackage{calc}      % table minipage width calculation for pandoc >= 2.11.1
    \usepackage[inline]{enumitem} % IRkernel/repr support (it uses the enumerate* environment)
    \usepackage[normalem]{ulem} % ulem is needed to support strikethroughs (\sout)
                                % normalem makes italics be italics, not underlines
    \usepackage{soul}      % strikethrough (\st) support for pandoc >= 3.0.0
    \usepackage{mathrsfs}
    

    
    % Colors for the hyperref package
    \definecolor{urlcolor}{rgb}{0,.145,.698}
    \definecolor{linkcolor}{rgb}{.71,0.21,0.01}
    \definecolor{citecolor}{rgb}{.12,.54,.11}

    % ANSI colors
    \definecolor{ansi-black}{HTML}{3E424D}
    \definecolor{ansi-black-intense}{HTML}{282C36}
    \definecolor{ansi-red}{HTML}{E75C58}
    \definecolor{ansi-red-intense}{HTML}{B22B31}
    \definecolor{ansi-green}{HTML}{00A250}
    \definecolor{ansi-green-intense}{HTML}{007427}
    \definecolor{ansi-yellow}{HTML}{DDB62B}
    \definecolor{ansi-yellow-intense}{HTML}{B27D12}
    \definecolor{ansi-blue}{HTML}{208FFB}
    \definecolor{ansi-blue-intense}{HTML}{0065CA}
    \definecolor{ansi-magenta}{HTML}{D160C4}
    \definecolor{ansi-magenta-intense}{HTML}{A03196}
    \definecolor{ansi-cyan}{HTML}{60C6C8}
    \definecolor{ansi-cyan-intense}{HTML}{258F8F}
    \definecolor{ansi-white}{HTML}{C5C1B4}
    \definecolor{ansi-white-intense}{HTML}{A1A6B2}
    \definecolor{ansi-default-inverse-fg}{HTML}{FFFFFF}
    \definecolor{ansi-default-inverse-bg}{HTML}{000000}

    % common color for the border for error outputs.
    \definecolor{outerrorbackground}{HTML}{FFDFDF}

    % commands and environments needed by pandoc snippets
    % extracted from the output of `pandoc -s`
    \providecommand{\tightlist}{%
      \setlength{\itemsep}{0pt}\setlength{\parskip}{0pt}}
    \DefineVerbatimEnvironment{Highlighting}{Verbatim}{commandchars=\\\{\}}
    % Add ',fontsize=\small' for more characters per line
    \newenvironment{Shaded}{}{}
    \newcommand{\KeywordTok}[1]{\textcolor[rgb]{0.00,0.44,0.13}{\textbf{{#1}}}}
    \newcommand{\DataTypeTok}[1]{\textcolor[rgb]{0.56,0.13,0.00}{{#1}}}
    \newcommand{\DecValTok}[1]{\textcolor[rgb]{0.25,0.63,0.44}{{#1}}}
    \newcommand{\BaseNTok}[1]{\textcolor[rgb]{0.25,0.63,0.44}{{#1}}}
    \newcommand{\FloatTok}[1]{\textcolor[rgb]{0.25,0.63,0.44}{{#1}}}
    \newcommand{\CharTok}[1]{\textcolor[rgb]{0.25,0.44,0.63}{{#1}}}
    \newcommand{\StringTok}[1]{\textcolor[rgb]{0.25,0.44,0.63}{{#1}}}
    \newcommand{\CommentTok}[1]{\textcolor[rgb]{0.38,0.63,0.69}{\textit{{#1}}}}
    \newcommand{\OtherTok}[1]{\textcolor[rgb]{0.00,0.44,0.13}{{#1}}}
    \newcommand{\AlertTok}[1]{\textcolor[rgb]{1.00,0.00,0.00}{\textbf{{#1}}}}
    \newcommand{\FunctionTok}[1]{\textcolor[rgb]{0.02,0.16,0.49}{{#1}}}
    \newcommand{\RegionMarkerTok}[1]{{#1}}
    \newcommand{\ErrorTok}[1]{\textcolor[rgb]{1.00,0.00,0.00}{\textbf{{#1}}}}
    \newcommand{\NormalTok}[1]{{#1}}

    % Additional commands for more recent versions of Pandoc
    \newcommand{\ConstantTok}[1]{\textcolor[rgb]{0.53,0.00,0.00}{{#1}}}
    \newcommand{\SpecialCharTok}[1]{\textcolor[rgb]{0.25,0.44,0.63}{{#1}}}
    \newcommand{\VerbatimStringTok}[1]{\textcolor[rgb]{0.25,0.44,0.63}{{#1}}}
    \newcommand{\SpecialStringTok}[1]{\textcolor[rgb]{0.73,0.40,0.53}{{#1}}}
    \newcommand{\ImportTok}[1]{{#1}}
    \newcommand{\DocumentationTok}[1]{\textcolor[rgb]{0.73,0.13,0.13}{\textit{{#1}}}}
    \newcommand{\AnnotationTok}[1]{\textcolor[rgb]{0.38,0.63,0.69}{\textbf{\textit{{#1}}}}}
    \newcommand{\CommentVarTok}[1]{\textcolor[rgb]{0.38,0.63,0.69}{\textbf{\textit{{#1}}}}}
    \newcommand{\VariableTok}[1]{\textcolor[rgb]{0.10,0.09,0.49}{{#1}}}
    \newcommand{\ControlFlowTok}[1]{\textcolor[rgb]{0.00,0.44,0.13}{\textbf{{#1}}}}
    \newcommand{\OperatorTok}[1]{\textcolor[rgb]{0.40,0.40,0.40}{{#1}}}
    \newcommand{\BuiltInTok}[1]{{#1}}
    \newcommand{\ExtensionTok}[1]{{#1}}
    \newcommand{\PreprocessorTok}[1]{\textcolor[rgb]{0.74,0.48,0.00}{{#1}}}
    \newcommand{\AttributeTok}[1]{\textcolor[rgb]{0.49,0.56,0.16}{{#1}}}
    \newcommand{\InformationTok}[1]{\textcolor[rgb]{0.38,0.63,0.69}{\textbf{\textit{{#1}}}}}
    \newcommand{\WarningTok}[1]{\textcolor[rgb]{0.38,0.63,0.69}{\textbf{\textit{{#1}}}}}
    \makeatletter
    \newsavebox\pandoc@box
    \newcommand*\pandocbounded[1]{%
      \sbox\pandoc@box{#1}%
      % scaling factors for width and height
      \Gscale@div\@tempa\textheight{\dimexpr\ht\pandoc@box+\dp\pandoc@box\relax}%
      \Gscale@div\@tempb\linewidth{\wd\pandoc@box}%
      % select the smaller of both
      \ifdim\@tempb\p@<\@tempa\p@
        \let\@tempa\@tempb
      \fi
      % scaling accordingly (\@tempa < 1)
      \ifdim\@tempa\p@<\p@
        \scalebox{\@tempa}{\usebox\pandoc@box}%
      % scaling not needed, use as it is
      \else
        \usebox{\pandoc@box}%
      \fi
    }
    \makeatother

    % Define a nice break command that doesn't care if a line doesn't already
    % exist.
    \def\br{\hspace*{\fill} \\* }
    % Math Jax compatibility definitions
    \def\gt{>}
    \def\lt{<}
    \let\Oldtex\TeX
    \let\Oldlatex\LaTeX
    \renewcommand{\TeX}{\textrm{\Oldtex}}
    \renewcommand{\LaTeX}{\textrm{\Oldlatex}}
    % Document parameters
    % Document title
    \title{Lab1}
    
    
    
    
    
    
    
% Pygments definitions
\makeatletter
\def\PY@reset{\let\PY@it=\relax \let\PY@bf=\relax%
    \let\PY@ul=\relax \let\PY@tc=\relax%
    \let\PY@bc=\relax \let\PY@ff=\relax}
\def\PY@tok#1{\csname PY@tok@#1\endcsname}
\def\PY@toks#1+{\ifx\relax#1\empty\else%
    \PY@tok{#1}\expandafter\PY@toks\fi}
\def\PY@do#1{\PY@bc{\PY@tc{\PY@ul{%
    \PY@it{\PY@bf{\PY@ff{#1}}}}}}}
\def\PY#1#2{\PY@reset\PY@toks#1+\relax+\PY@do{#2}}

\@namedef{PY@tok@w}{\def\PY@tc##1{\textcolor[rgb]{0.73,0.73,0.73}{##1}}}
\@namedef{PY@tok@c}{\let\PY@it=\textit\def\PY@tc##1{\textcolor[rgb]{0.24,0.48,0.48}{##1}}}
\@namedef{PY@tok@cp}{\def\PY@tc##1{\textcolor[rgb]{0.61,0.40,0.00}{##1}}}
\@namedef{PY@tok@k}{\let\PY@bf=\textbf\def\PY@tc##1{\textcolor[rgb]{0.00,0.50,0.00}{##1}}}
\@namedef{PY@tok@kp}{\def\PY@tc##1{\textcolor[rgb]{0.00,0.50,0.00}{##1}}}
\@namedef{PY@tok@kt}{\def\PY@tc##1{\textcolor[rgb]{0.69,0.00,0.25}{##1}}}
\@namedef{PY@tok@o}{\def\PY@tc##1{\textcolor[rgb]{0.40,0.40,0.40}{##1}}}
\@namedef{PY@tok@ow}{\let\PY@bf=\textbf\def\PY@tc##1{\textcolor[rgb]{0.67,0.13,1.00}{##1}}}
\@namedef{PY@tok@nb}{\def\PY@tc##1{\textcolor[rgb]{0.00,0.50,0.00}{##1}}}
\@namedef{PY@tok@nf}{\def\PY@tc##1{\textcolor[rgb]{0.00,0.00,1.00}{##1}}}
\@namedef{PY@tok@nc}{\let\PY@bf=\textbf\def\PY@tc##1{\textcolor[rgb]{0.00,0.00,1.00}{##1}}}
\@namedef{PY@tok@nn}{\let\PY@bf=\textbf\def\PY@tc##1{\textcolor[rgb]{0.00,0.00,1.00}{##1}}}
\@namedef{PY@tok@ne}{\let\PY@bf=\textbf\def\PY@tc##1{\textcolor[rgb]{0.80,0.25,0.22}{##1}}}
\@namedef{PY@tok@nv}{\def\PY@tc##1{\textcolor[rgb]{0.10,0.09,0.49}{##1}}}
\@namedef{PY@tok@no}{\def\PY@tc##1{\textcolor[rgb]{0.53,0.00,0.00}{##1}}}
\@namedef{PY@tok@nl}{\def\PY@tc##1{\textcolor[rgb]{0.46,0.46,0.00}{##1}}}
\@namedef{PY@tok@ni}{\let\PY@bf=\textbf\def\PY@tc##1{\textcolor[rgb]{0.44,0.44,0.44}{##1}}}
\@namedef{PY@tok@na}{\def\PY@tc##1{\textcolor[rgb]{0.41,0.47,0.13}{##1}}}
\@namedef{PY@tok@nt}{\let\PY@bf=\textbf\def\PY@tc##1{\textcolor[rgb]{0.00,0.50,0.00}{##1}}}
\@namedef{PY@tok@nd}{\def\PY@tc##1{\textcolor[rgb]{0.67,0.13,1.00}{##1}}}
\@namedef{PY@tok@s}{\def\PY@tc##1{\textcolor[rgb]{0.73,0.13,0.13}{##1}}}
\@namedef{PY@tok@sd}{\let\PY@it=\textit\def\PY@tc##1{\textcolor[rgb]{0.73,0.13,0.13}{##1}}}
\@namedef{PY@tok@si}{\let\PY@bf=\textbf\def\PY@tc##1{\textcolor[rgb]{0.64,0.35,0.47}{##1}}}
\@namedef{PY@tok@se}{\let\PY@bf=\textbf\def\PY@tc##1{\textcolor[rgb]{0.67,0.36,0.12}{##1}}}
\@namedef{PY@tok@sr}{\def\PY@tc##1{\textcolor[rgb]{0.64,0.35,0.47}{##1}}}
\@namedef{PY@tok@ss}{\def\PY@tc##1{\textcolor[rgb]{0.10,0.09,0.49}{##1}}}
\@namedef{PY@tok@sx}{\def\PY@tc##1{\textcolor[rgb]{0.00,0.50,0.00}{##1}}}
\@namedef{PY@tok@m}{\def\PY@tc##1{\textcolor[rgb]{0.40,0.40,0.40}{##1}}}
\@namedef{PY@tok@gh}{\let\PY@bf=\textbf\def\PY@tc##1{\textcolor[rgb]{0.00,0.00,0.50}{##1}}}
\@namedef{PY@tok@gu}{\let\PY@bf=\textbf\def\PY@tc##1{\textcolor[rgb]{0.50,0.00,0.50}{##1}}}
\@namedef{PY@tok@gd}{\def\PY@tc##1{\textcolor[rgb]{0.63,0.00,0.00}{##1}}}
\@namedef{PY@tok@gi}{\def\PY@tc##1{\textcolor[rgb]{0.00,0.52,0.00}{##1}}}
\@namedef{PY@tok@gr}{\def\PY@tc##1{\textcolor[rgb]{0.89,0.00,0.00}{##1}}}
\@namedef{PY@tok@ge}{\let\PY@it=\textit}
\@namedef{PY@tok@gs}{\let\PY@bf=\textbf}
\@namedef{PY@tok@ges}{\let\PY@bf=\textbf\let\PY@it=\textit}
\@namedef{PY@tok@gp}{\let\PY@bf=\textbf\def\PY@tc##1{\textcolor[rgb]{0.00,0.00,0.50}{##1}}}
\@namedef{PY@tok@go}{\def\PY@tc##1{\textcolor[rgb]{0.44,0.44,0.44}{##1}}}
\@namedef{PY@tok@gt}{\def\PY@tc##1{\textcolor[rgb]{0.00,0.27,0.87}{##1}}}
\@namedef{PY@tok@err}{\def\PY@bc##1{{\setlength{\fboxsep}{\string -\fboxrule}\fcolorbox[rgb]{1.00,0.00,0.00}{1,1,1}{\strut ##1}}}}
\@namedef{PY@tok@kc}{\let\PY@bf=\textbf\def\PY@tc##1{\textcolor[rgb]{0.00,0.50,0.00}{##1}}}
\@namedef{PY@tok@kd}{\let\PY@bf=\textbf\def\PY@tc##1{\textcolor[rgb]{0.00,0.50,0.00}{##1}}}
\@namedef{PY@tok@kn}{\let\PY@bf=\textbf\def\PY@tc##1{\textcolor[rgb]{0.00,0.50,0.00}{##1}}}
\@namedef{PY@tok@kr}{\let\PY@bf=\textbf\def\PY@tc##1{\textcolor[rgb]{0.00,0.50,0.00}{##1}}}
\@namedef{PY@tok@bp}{\def\PY@tc##1{\textcolor[rgb]{0.00,0.50,0.00}{##1}}}
\@namedef{PY@tok@fm}{\def\PY@tc##1{\textcolor[rgb]{0.00,0.00,1.00}{##1}}}
\@namedef{PY@tok@vc}{\def\PY@tc##1{\textcolor[rgb]{0.10,0.09,0.49}{##1}}}
\@namedef{PY@tok@vg}{\def\PY@tc##1{\textcolor[rgb]{0.10,0.09,0.49}{##1}}}
\@namedef{PY@tok@vi}{\def\PY@tc##1{\textcolor[rgb]{0.10,0.09,0.49}{##1}}}
\@namedef{PY@tok@vm}{\def\PY@tc##1{\textcolor[rgb]{0.10,0.09,0.49}{##1}}}
\@namedef{PY@tok@sa}{\def\PY@tc##1{\textcolor[rgb]{0.73,0.13,0.13}{##1}}}
\@namedef{PY@tok@sb}{\def\PY@tc##1{\textcolor[rgb]{0.73,0.13,0.13}{##1}}}
\@namedef{PY@tok@sc}{\def\PY@tc##1{\textcolor[rgb]{0.73,0.13,0.13}{##1}}}
\@namedef{PY@tok@dl}{\def\PY@tc##1{\textcolor[rgb]{0.73,0.13,0.13}{##1}}}
\@namedef{PY@tok@s2}{\def\PY@tc##1{\textcolor[rgb]{0.73,0.13,0.13}{##1}}}
\@namedef{PY@tok@sh}{\def\PY@tc##1{\textcolor[rgb]{0.73,0.13,0.13}{##1}}}
\@namedef{PY@tok@s1}{\def\PY@tc##1{\textcolor[rgb]{0.73,0.13,0.13}{##1}}}
\@namedef{PY@tok@mb}{\def\PY@tc##1{\textcolor[rgb]{0.40,0.40,0.40}{##1}}}
\@namedef{PY@tok@mf}{\def\PY@tc##1{\textcolor[rgb]{0.40,0.40,0.40}{##1}}}
\@namedef{PY@tok@mh}{\def\PY@tc##1{\textcolor[rgb]{0.40,0.40,0.40}{##1}}}
\@namedef{PY@tok@mi}{\def\PY@tc##1{\textcolor[rgb]{0.40,0.40,0.40}{##1}}}
\@namedef{PY@tok@il}{\def\PY@tc##1{\textcolor[rgb]{0.40,0.40,0.40}{##1}}}
\@namedef{PY@tok@mo}{\def\PY@tc##1{\textcolor[rgb]{0.40,0.40,0.40}{##1}}}
\@namedef{PY@tok@ch}{\let\PY@it=\textit\def\PY@tc##1{\textcolor[rgb]{0.24,0.48,0.48}{##1}}}
\@namedef{PY@tok@cm}{\let\PY@it=\textit\def\PY@tc##1{\textcolor[rgb]{0.24,0.48,0.48}{##1}}}
\@namedef{PY@tok@cpf}{\let\PY@it=\textit\def\PY@tc##1{\textcolor[rgb]{0.24,0.48,0.48}{##1}}}
\@namedef{PY@tok@c1}{\let\PY@it=\textit\def\PY@tc##1{\textcolor[rgb]{0.24,0.48,0.48}{##1}}}
\@namedef{PY@tok@cs}{\let\PY@it=\textit\def\PY@tc##1{\textcolor[rgb]{0.24,0.48,0.48}{##1}}}

\def\PYZbs{\char`\\}
\def\PYZus{\char`\_}
\def\PYZob{\char`\{}
\def\PYZcb{\char`\}}
\def\PYZca{\char`\^}
\def\PYZam{\char`\&}
\def\PYZlt{\char`\<}
\def\PYZgt{\char`\>}
\def\PYZsh{\char`\#}
\def\PYZpc{\char`\%}
\def\PYZdl{\char`\$}
\def\PYZhy{\char`\-}
\def\PYZsq{\char`\'}
\def\PYZdq{\char`\"}
\def\PYZti{\char`\~}
% for compatibility with earlier versions
\def\PYZat{@}
\def\PYZlb{[}
\def\PYZrb{]}
\makeatother


    % For linebreaks inside Verbatim environment from package fancyvrb.
    \makeatletter
        \newbox\Wrappedcontinuationbox
        \newbox\Wrappedvisiblespacebox
        \newcommand*\Wrappedvisiblespace {\textcolor{red}{\textvisiblespace}}
        \newcommand*\Wrappedcontinuationsymbol {\textcolor{red}{\llap{\tiny$\m@th\hookrightarrow$}}}
        \newcommand*\Wrappedcontinuationindent {3ex }
        \newcommand*\Wrappedafterbreak {\kern\Wrappedcontinuationindent\copy\Wrappedcontinuationbox}
        % Take advantage of the already applied Pygments mark-up to insert
        % potential linebreaks for TeX processing.
        %        {, <, #, %, $, ' and ": go to next line.
        %        _, }, ^, &, >, - and ~: stay at end of broken line.
        % Use of \textquotesingle for straight quote.
        \newcommand*\Wrappedbreaksatspecials {%
            \def\PYGZus{\discretionary{\char`\_}{\Wrappedafterbreak}{\char`\_}}%
            \def\PYGZob{\discretionary{}{\Wrappedafterbreak\char`\{}{\char`\{}}%
            \def\PYGZcb{\discretionary{\char`\}}{\Wrappedafterbreak}{\char`\}}}%
            \def\PYGZca{\discretionary{\char`\^}{\Wrappedafterbreak}{\char`\^}}%
            \def\PYGZam{\discretionary{\char`\&}{\Wrappedafterbreak}{\char`\&}}%
            \def\PYGZlt{\discretionary{}{\Wrappedafterbreak\char`\<}{\char`\<}}%
            \def\PYGZgt{\discretionary{\char`\>}{\Wrappedafterbreak}{\char`\>}}%
            \def\PYGZsh{\discretionary{}{\Wrappedafterbreak\char`\#}{\char`\#}}%
            \def\PYGZpc{\discretionary{}{\Wrappedafterbreak\char`\%}{\char`\%}}%
            \def\PYGZdl{\discretionary{}{\Wrappedafterbreak\char`\$}{\char`\$}}%
            \def\PYGZhy{\discretionary{\char`\-}{\Wrappedafterbreak}{\char`\-}}%
            \def\PYGZsq{\discretionary{}{\Wrappedafterbreak\textquotesingle}{\textquotesingle}}%
            \def\PYGZdq{\discretionary{}{\Wrappedafterbreak\char`\"}{\char`\"}}%
            \def\PYGZti{\discretionary{\char`\~}{\Wrappedafterbreak}{\char`\~}}%
        }
        % Some characters . , ; ? ! / are not pygmentized.
        % This macro makes them "active" and they will insert potential linebreaks
        \newcommand*\Wrappedbreaksatpunct {%
            \lccode`\~`\.\lowercase{\def~}{\discretionary{\hbox{\char`\.}}{\Wrappedafterbreak}{\hbox{\char`\.}}}%
            \lccode`\~`\,\lowercase{\def~}{\discretionary{\hbox{\char`\,}}{\Wrappedafterbreak}{\hbox{\char`\,}}}%
            \lccode`\~`\;\lowercase{\def~}{\discretionary{\hbox{\char`\;}}{\Wrappedafterbreak}{\hbox{\char`\;}}}%
            \lccode`\~`\:\lowercase{\def~}{\discretionary{\hbox{\char`\:}}{\Wrappedafterbreak}{\hbox{\char`\:}}}%
            \lccode`\~`\?\lowercase{\def~}{\discretionary{\hbox{\char`\?}}{\Wrappedafterbreak}{\hbox{\char`\?}}}%
            \lccode`\~`\!\lowercase{\def~}{\discretionary{\hbox{\char`\!}}{\Wrappedafterbreak}{\hbox{\char`\!}}}%
            \lccode`\~`\/\lowercase{\def~}{\discretionary{\hbox{\char`\/}}{\Wrappedafterbreak}{\hbox{\char`\/}}}%
            \catcode`\.\active
            \catcode`\,\active
            \catcode`\;\active
            \catcode`\:\active
            \catcode`\?\active
            \catcode`\!\active
            \catcode`\/\active
            \lccode`\~`\~
        }
    \makeatother

    \let\OriginalVerbatim=\Verbatim
    \makeatletter
    \renewcommand{\Verbatim}[1][1]{%
        %\parskip\z@skip
        \sbox\Wrappedcontinuationbox {\Wrappedcontinuationsymbol}%
        \sbox\Wrappedvisiblespacebox {\FV@SetupFont\Wrappedvisiblespace}%
        \def\FancyVerbFormatLine ##1{\hsize\linewidth
            \vtop{\raggedright\hyphenpenalty\z@\exhyphenpenalty\z@
                \doublehyphendemerits\z@\finalhyphendemerits\z@
                \strut ##1\strut}%
        }%
        % If the linebreak is at a space, the latter will be displayed as visible
        % space at end of first line, and a continuation symbol starts next line.
        % Stretch/shrink are however usually zero for typewriter font.
        \def\FV@Space {%
            \nobreak\hskip\z@ plus\fontdimen3\font minus\fontdimen4\font
            \discretionary{\copy\Wrappedvisiblespacebox}{\Wrappedafterbreak}
            {\kern\fontdimen2\font}%
        }%

        % Allow breaks at special characters using \PYG... macros.
        \Wrappedbreaksatspecials
        % Breaks at punctuation characters . , ; ? ! and / need catcode=\active
        \OriginalVerbatim[#1,codes*=\Wrappedbreaksatpunct]%
    }
    \makeatother

    % Exact colors from NB
    \definecolor{incolor}{HTML}{303F9F}
    \definecolor{outcolor}{HTML}{D84315}
    \definecolor{cellborder}{HTML}{CFCFCF}
    \definecolor{cellbackground}{HTML}{F7F7F7}

    % prompt
    \makeatletter
    \newcommand{\boxspacing}{\kern\kvtcb@left@rule\kern\kvtcb@boxsep}
    \makeatother
    \newcommand{\prompt}[4]{
        {\ttfamily\llap{{\color{#2}[#3]:\hspace{3pt}#4}}\vspace{-\baselineskip}}
    }
    

    
    % Prevent overflowing lines due to hard-to-break entities
    \sloppy
    % Setup hyperref package
    \hypersetup{
      breaklinks=true,  % so long urls are correctly broken across lines
      colorlinks=true,
      urlcolor=urlcolor,
      linkcolor=linkcolor,
      citecolor=citecolor,
      }
    % Slightly bigger margins than the latex defaults
    
    \geometry{verbose,tmargin=1in,bmargin=1in,lmargin=1in,rmargin=1in}
    
    

\begin{document}
    
    \maketitle
    
    

    
    \begin{tcolorbox}[breakable, size=fbox, boxrule=1pt, pad at break*=1mm,colback=cellbackground, colframe=cellborder]
\prompt{In}{incolor}{1}{\boxspacing}
\begin{Verbatim}[commandchars=\\\{\}]
\PY{k+kn}{import}\PY{+w}{ }\PY{n+nn}{torch}
\PY{k+kn}{import}\PY{+w}{ }\PY{n+nn}{time}
\PY{k+kn}{import}\PY{+w}{ }\PY{n+nn}{numpy}\PY{+w}{ }\PY{k}{as}\PY{+w}{ }\PY{n+nn}{np}
\PY{k+kn}{import}\PY{+w}{ }\PY{n+nn}{matplotlib}\PY{n+nn}{.}\PY{n+nn}{pyplot}\PY{+w}{ }\PY{k}{as}\PY{+w}{ }\PY{n+nn}{plt}

\PY{n+nb}{print}\PY{p}{(}\PY{l+s+sa}{f}\PY{l+s+s2}{\PYZdq{}}\PY{l+s+s2}{PyTorch version: }\PY{l+s+si}{\PYZob{}}\PY{n}{torch}\PY{o}{.}\PY{n}{\PYZus{}\PYZus{}version\PYZus{}\PYZus{}}\PY{l+s+si}{\PYZcb{}}\PY{l+s+s2}{\PYZdq{}}\PY{p}{)}
\PY{n+nb}{print}\PY{p}{(}\PY{l+s+sa}{f}\PY{l+s+s2}{\PYZdq{}}\PY{l+s+s2}{CUDA available: }\PY{l+s+si}{\PYZob{}}\PY{n}{torch}\PY{o}{.}\PY{n}{cuda}\PY{o}{.}\PY{n}{is\PYZus{}available}\PY{p}{(}\PY{p}{)}\PY{l+s+si}{\PYZcb{}}\PY{l+s+s2}{\PYZdq{}}\PY{p}{)}
\PY{n+nb}{print}\PY{p}{(}\PY{l+s+sa}{f}\PY{l+s+s2}{\PYZdq{}}\PY{l+s+s2}{MPS available: }\PY{l+s+si}{\PYZob{}}\PY{n}{torch}\PY{o}{.}\PY{n}{backends}\PY{o}{.}\PY{n}{mps}\PY{o}{.}\PY{n}{is\PYZus{}available}\PY{p}{(}\PY{p}{)}\PY{l+s+si}{\PYZcb{}}\PY{l+s+s2}{\PYZdq{}}\PY{p}{)}

\PY{c+c1}{\PYZsh{} Determine the best available device}
\PY{k}{if} \PY{n}{torch}\PY{o}{.}\PY{n}{cuda}\PY{o}{.}\PY{n}{is\PYZus{}available}\PY{p}{(}\PY{p}{)}\PY{p}{:}
    \PY{n}{device} \PY{o}{=} \PY{n}{torch}\PY{o}{.}\PY{n}{device}\PY{p}{(}\PY{l+s+s2}{\PYZdq{}}\PY{l+s+s2}{cuda}\PY{l+s+s2}{\PYZdq{}}\PY{p}{)}
    \PY{n}{device\PYZus{}name} \PY{o}{=} \PY{n}{torch}\PY{o}{.}\PY{n}{cuda}\PY{o}{.}\PY{n}{get\PYZus{}device\PYZus{}name}\PY{p}{(}\PY{l+m+mi}{0}\PY{p}{)}
    \PY{n+nb}{print}\PY{p}{(}\PY{l+s+sa}{f}\PY{l+s+s2}{\PYZdq{}}\PY{l+s+s2}{Using CUDA device: }\PY{l+s+si}{\PYZob{}}\PY{n}{device\PYZus{}name}\PY{l+s+si}{\PYZcb{}}\PY{l+s+s2}{\PYZdq{}}\PY{p}{)}
\PY{k}{elif} \PY{n}{torch}\PY{o}{.}\PY{n}{backends}\PY{o}{.}\PY{n}{mps}\PY{o}{.}\PY{n}{is\PYZus{}available}\PY{p}{(}\PY{p}{)}\PY{p}{:}
    \PY{n}{device} \PY{o}{=} \PY{n}{torch}\PY{o}{.}\PY{n}{device}\PY{p}{(}\PY{l+s+s2}{\PYZdq{}}\PY{l+s+s2}{mps}\PY{l+s+s2}{\PYZdq{}}\PY{p}{)}
    \PY{n}{device\PYZus{}name} \PY{o}{=} \PY{l+s+s2}{\PYZdq{}}\PY{l+s+s2}{Apple Silicon (MPS)}\PY{l+s+s2}{\PYZdq{}}
    \PY{n+nb}{print}\PY{p}{(}\PY{l+s+sa}{f}\PY{l+s+s2}{\PYZdq{}}\PY{l+s+s2}{Using MPS device: }\PY{l+s+si}{\PYZob{}}\PY{n}{device\PYZus{}name}\PY{l+s+si}{\PYZcb{}}\PY{l+s+s2}{\PYZdq{}}\PY{p}{)}
\PY{k}{else}\PY{p}{:}
    \PY{n}{device} \PY{o}{=} \PY{n}{torch}\PY{o}{.}\PY{n}{device}\PY{p}{(}\PY{l+s+s2}{\PYZdq{}}\PY{l+s+s2}{cpu}\PY{l+s+s2}{\PYZdq{}}\PY{p}{)}
    \PY{n}{device\PYZus{}name} \PY{o}{=} \PY{l+s+s2}{\PYZdq{}}\PY{l+s+s2}{CPU only}\PY{l+s+s2}{\PYZdq{}}
    \PY{n+nb}{print}\PY{p}{(}\PY{l+s+s2}{\PYZdq{}}\PY{l+s+s2}{No GPU available, using CPU}\PY{l+s+s2}{\PYZdq{}}\PY{p}{)}
    \PY{n+nb}{print}\PY{p}{(}\PY{l+s+s2}{\PYZdq{}}\PY{l+s+s2}{WARNING: Please use Google Colab for this assignment!}\PY{l+s+s2}{\PYZdq{}}\PY{p}{)}
\end{Verbatim}
\end{tcolorbox}

    \begin{Verbatim}[commandchars=\\\{\}]
PyTorch version: 2.10.0+cu130
CUDA available: True
MPS available: False
Using CUDA device: NVIDIA GeForce RTX 4060 Laptop GPU
    \end{Verbatim}

    \begin{tcolorbox}[breakable, size=fbox, boxrule=1pt, pad at break*=1mm,colback=cellbackground, colframe=cellborder]
\prompt{In}{incolor}{2}{\boxspacing}
\begin{Verbatim}[commandchars=\\\{\}]
\PY{k}{def}\PY{+w}{ }\PY{n+nf}{create\PYZus{}matrices}\PY{p}{(}\PY{n}{size}\PY{p}{,} \PY{n}{device}\PY{p}{)}\PY{p}{:}
\PY{+w}{    }\PY{l+s+sd}{\PYZdq{}\PYZdq{}\PYZdq{}}
\PY{l+s+sd}{    Create two random matrices of shape (size, size).}

\PY{l+s+sd}{    Args:}
\PY{l+s+sd}{        size: The dimension of the square matrices.}
\PY{l+s+sd}{        device: The device to create tensors on (\PYZsq{}cpu\PYZsq{}, \PYZsq{}cuda\PYZsq{}, or \PYZsq{}mps\PYZsq{}).}

\PY{l+s+sd}{    Returns:}
\PY{l+s+sd}{        Two PyTorch tensors on the specified device.}
\PY{l+s+sd}{    \PYZdq{}\PYZdq{}\PYZdq{}}
    \PY{n}{A} \PY{o}{=} \PY{n}{torch}\PY{o}{.}\PY{n}{randn}\PY{p}{(}\PY{n}{size}\PY{p}{,} \PY{n}{size}\PY{p}{,} \PY{n}{dtype}\PY{o}{=}\PY{n}{torch}\PY{o}{.}\PY{n}{float32}\PY{p}{,} \PY{n}{device}\PY{o}{=}\PY{n}{device}\PY{p}{)}
    \PY{n}{B} \PY{o}{=} \PY{n}{torch}\PY{o}{.}\PY{n}{randn}\PY{p}{(}\PY{n}{size}\PY{p}{,} \PY{n}{size}\PY{p}{,} \PY{n}{dtype}\PY{o}{=}\PY{n}{torch}\PY{o}{.}\PY{n}{float32}\PY{p}{,} \PY{n}{device}\PY{o}{=}\PY{n}{device}\PY{p}{)}
    
    \PY{k}{return} \PY{n}{A}\PY{p}{,} \PY{n}{B}
\end{Verbatim}
\end{tcolorbox}

    \begin{tcolorbox}[breakable, size=fbox, boxrule=1pt, pad at break*=1mm,colback=cellbackground, colframe=cellborder]
\prompt{In}{incolor}{3}{\boxspacing}
\begin{Verbatim}[commandchars=\\\{\}]
\PY{n}{size} \PY{o}{=} \PY{l+m+mi}{1000}

\PY{n}{A\PYZus{}cpu}\PY{p}{,} \PY{n}{B\PYZus{}cpu} \PY{o}{=} \PY{n}{create\PYZus{}matrices}\PY{p}{(}\PY{n}{size}\PY{p}{,} \PY{n}{torch}\PY{o}{.}\PY{n}{device}\PY{p}{(}\PY{l+s+s2}{\PYZdq{}}\PY{l+s+s2}{cpu}\PY{l+s+s2}{\PYZdq{}}\PY{p}{)}\PY{p}{)}

\PY{n+nb}{print}\PY{p}{(}\PY{l+s+sa}{f}\PY{l+s+s2}{\PYZdq{}}\PY{l+s+s2}{Matrix A shape: }\PY{l+s+si}{\PYZob{}}\PY{n}{A\PYZus{}cpu}\PY{o}{.}\PY{n}{shape}\PY{l+s+si}{\PYZcb{}}\PY{l+s+s2}{\PYZdq{}}\PY{p}{)}
\PY{n+nb}{print}\PY{p}{(}\PY{l+s+sa}{f}\PY{l+s+s2}{\PYZdq{}}\PY{l+s+s2}{Matrix A device: }\PY{l+s+si}{\PYZob{}}\PY{n}{A\PYZus{}cpu}\PY{o}{.}\PY{n}{device}\PY{l+s+si}{\PYZcb{}}\PY{l+s+s2}{\PYZdq{}}\PY{p}{)}
\PY{n+nb}{print}\PY{p}{(}\PY{l+s+sa}{f}\PY{l+s+s2}{\PYZdq{}}\PY{l+s+s2}{Matrix A dtype: }\PY{l+s+si}{\PYZob{}}\PY{n}{A\PYZus{}cpu}\PY{o}{.}\PY{n}{dtype}\PY{l+s+si}{\PYZcb{}}\PY{l+s+s2}{\PYZdq{}}\PY{p}{)}
\end{Verbatim}
\end{tcolorbox}

    \begin{Verbatim}[commandchars=\\\{\}]
Matrix A shape: torch.Size([1000, 1000])
Matrix A device: cpu
Matrix A dtype: torch.float32
    \end{Verbatim}

    \begin{tcolorbox}[breakable, size=fbox, boxrule=1pt, pad at break*=1mm,colback=cellbackground, colframe=cellborder]
\prompt{In}{incolor}{4}{\boxspacing}
\begin{Verbatim}[commandchars=\\\{\}]
\PY{k}{def}\PY{+w}{ }\PY{n+nf}{benchmark\PYZus{}matmul}\PY{p}{(}\PY{n}{A}\PY{p}{,} \PY{n}{B}\PY{p}{)}\PY{p}{:}
\PY{+w}{    }\PY{l+s+sd}{\PYZdq{}\PYZdq{}\PYZdq{}}
\PY{l+s+sd}{    Benchmark matrix multiplication with proper synchronization.}

\PY{l+s+sd}{    Args:}
\PY{l+s+sd}{        A, B: Input matrices}

\PY{l+s+sd}{    Returns:}
\PY{l+s+sd}{        Execution time in milliseconds}
\PY{l+s+sd}{    \PYZdq{}\PYZdq{}\PYZdq{}}
    \PY{n}{device} \PY{o}{=} \PY{n}{A}\PY{o}{.}\PY{n}{device}

    \PY{c+c1}{\PYZsh{} Synchronize before timing}
    \PY{k}{if} \PY{n}{device}\PY{o}{.}\PY{n}{type} \PY{o}{==} \PY{l+s+s2}{\PYZdq{}}\PY{l+s+s2}{cuda}\PY{l+s+s2}{\PYZdq{}}\PY{p}{:}
        \PY{n}{torch}\PY{o}{.}\PY{n}{cuda}\PY{o}{.}\PY{n}{synchronize}\PY{p}{(}\PY{p}{)}
    \PY{k}{elif} \PY{n}{device}\PY{o}{.}\PY{n}{type} \PY{o}{==} \PY{l+s+s2}{\PYZdq{}}\PY{l+s+s2}{mps}\PY{l+s+s2}{\PYZdq{}}\PY{p}{:}
        \PY{n}{torch}\PY{o}{.}\PY{n}{mps}\PY{o}{.}\PY{n}{synchronize}\PY{p}{(}\PY{p}{)}

    \PY{n}{start} \PY{o}{=} \PY{n}{time}\PY{o}{.}\PY{n}{perf\PYZus{}counter}\PY{p}{(}\PY{p}{)}

    \PY{n}{C} \PY{o}{=} \PY{n}{torch}\PY{o}{.}\PY{n}{mm}\PY{p}{(}\PY{n}{A}\PY{p}{,} \PY{n}{B}\PY{p}{)}

    \PY{c+c1}{\PYZsh{} Synchronize after operation}
    \PY{k}{if} \PY{n}{device}\PY{o}{.}\PY{n}{type} \PY{o}{==} \PY{l+s+s2}{\PYZdq{}}\PY{l+s+s2}{cuda}\PY{l+s+s2}{\PYZdq{}}\PY{p}{:}
        \PY{n}{torch}\PY{o}{.}\PY{n}{cuda}\PY{o}{.}\PY{n}{synchronize}\PY{p}{(}\PY{p}{)}
    \PY{k}{elif} \PY{n}{device}\PY{o}{.}\PY{n}{type} \PY{o}{==} \PY{l+s+s2}{\PYZdq{}}\PY{l+s+s2}{mps}\PY{l+s+s2}{\PYZdq{}}\PY{p}{:}
        \PY{n}{torch}\PY{o}{.}\PY{n}{mps}\PY{o}{.}\PY{n}{synchronize}\PY{p}{(}\PY{p}{)}

    \PY{n}{end} \PY{o}{=} \PY{n}{time}\PY{o}{.}\PY{n}{perf\PYZus{}counter}\PY{p}{(}\PY{p}{)}

    \PY{k}{return} \PY{p}{(}\PY{n}{end} \PY{o}{\PYZhy{}} \PY{n}{start}\PY{p}{)} \PY{o}{*} \PY{l+m+mi}{1000}
\end{Verbatim}
\end{tcolorbox}

    \begin{tcolorbox}[breakable, size=fbox, boxrule=1pt, pad at break*=1mm,colback=cellbackground, colframe=cellborder]
\prompt{In}{incolor}{5}{\boxspacing}
\begin{Verbatim}[commandchars=\\\{\}]
\PY{c+c1}{\PYZsh{} Matrix sizes to test}
\PY{n}{sizes} \PY{o}{=} \PY{p}{[}\PY{l+m+mi}{512}\PY{p}{,} \PY{l+m+mi}{1024}\PY{p}{,} \PY{l+m+mi}{2048}\PY{p}{,} \PY{l+m+mi}{4096}\PY{p}{,} \PY{l+m+mi}{8192}\PY{p}{]}

\PY{c+c1}{\PYZsh{} Store results}
\PY{n}{results} \PY{o}{=} \PY{p}{\PYZob{}}
    \PY{l+s+s2}{\PYZdq{}}\PY{l+s+s2}{sizes}\PY{l+s+s2}{\PYZdq{}}\PY{p}{:} \PY{n}{sizes}\PY{p}{,}
    \PY{l+s+s2}{\PYZdq{}}\PY{l+s+s2}{cpu\PYZus{}times}\PY{l+s+s2}{\PYZdq{}}\PY{p}{:} \PY{p}{[}\PY{p}{]}\PY{p}{,}
    \PY{l+s+s2}{\PYZdq{}}\PY{l+s+s2}{gpu\PYZus{}times}\PY{l+s+s2}{\PYZdq{}}\PY{p}{:} \PY{p}{[}\PY{p}{]}\PY{p}{,}
    \PY{l+s+s2}{\PYZdq{}}\PY{l+s+s2}{speedups}\PY{l+s+s2}{\PYZdq{}}\PY{p}{:} \PY{p}{[}\PY{p}{]}
\PY{p}{\PYZcb{}}

\PY{n+nb}{print}\PY{p}{(}\PY{l+s+s2}{\PYZdq{}}\PY{l+s+s2}{=}\PY{l+s+s2}{\PYZdq{}} \PY{o}{*} \PY{l+m+mi}{60}\PY{p}{)}
\PY{n+nb}{print}\PY{p}{(}\PY{l+s+sa}{f}\PY{l+s+s2}{\PYZdq{}}\PY{l+s+si}{\PYZob{}}\PY{l+s+s1}{\PYZsq{}}\PY{l+s+s1}{Size}\PY{l+s+s1}{\PYZsq{}}\PY{l+s+si}{:}\PY{l+s+s2}{\PYZgt{}8}\PY{l+s+si}{\PYZcb{}}\PY{l+s+s2}{ | }\PY{l+s+si}{\PYZob{}}\PY{l+s+s1}{\PYZsq{}}\PY{l+s+s1}{CPU (ms)}\PY{l+s+s1}{\PYZsq{}}\PY{l+s+si}{:}\PY{l+s+s2}{\PYZgt{}12}\PY{l+s+si}{\PYZcb{}}\PY{l+s+s2}{ | }\PY{l+s+si}{\PYZob{}}\PY{l+s+s1}{\PYZsq{}}\PY{l+s+s1}{GPU (ms)}\PY{l+s+s1}{\PYZsq{}}\PY{l+s+si}{:}\PY{l+s+s2}{\PYZgt{}12}\PY{l+s+si}{\PYZcb{}}\PY{l+s+s2}{ | }\PY{l+s+si}{\PYZob{}}\PY{l+s+s1}{\PYZsq{}}\PY{l+s+s1}{Speedup}\PY{l+s+s1}{\PYZsq{}}\PY{l+s+si}{:}\PY{l+s+s2}{\PYZgt{}10}\PY{l+s+si}{\PYZcb{}}\PY{l+s+s2}{\PYZdq{}}\PY{p}{)}
\PY{n+nb}{print}\PY{p}{(}\PY{l+s+s2}{\PYZdq{}}\PY{l+s+s2}{=}\PY{l+s+s2}{\PYZdq{}} \PY{o}{*} \PY{l+m+mi}{60}\PY{p}{)}

\PY{k}{for} \PY{n}{size} \PY{o+ow}{in} \PY{n}{sizes}\PY{p}{:}
    \PY{c+c1}{\PYZsh{} CPU benchmark}
    \PY{n}{A\PYZus{}cpu}\PY{p}{,} \PY{n}{B\PYZus{}cpu} \PY{o}{=} \PY{n}{create\PYZus{}matrices}\PY{p}{(}\PY{n}{size}\PY{p}{,} \PY{n}{torch}\PY{o}{.}\PY{n}{device}\PY{p}{(}\PY{l+s+s2}{\PYZdq{}}\PY{l+s+s2}{cpu}\PY{l+s+s2}{\PYZdq{}}\PY{p}{)}\PY{p}{)}
    \PY{n}{cpu\PYZus{}time} \PY{o}{=} \PY{n}{benchmark\PYZus{}matmul}\PY{p}{(}\PY{n}{A\PYZus{}cpu}\PY{p}{,} \PY{n}{B\PYZus{}cpu}\PY{p}{)}

    \PY{c+c1}{\PYZsh{} GPU benchmark}
    \PY{n}{A\PYZus{}gpu}\PY{p}{,} \PY{n}{B\PYZus{}gpu} \PY{o}{=} \PY{n}{create\PYZus{}matrices}\PY{p}{(}\PY{n}{size}\PY{p}{,} \PY{n}{device}\PY{p}{)}
    \PY{n}{gpu\PYZus{}time} \PY{o}{=} \PY{n}{benchmark\PYZus{}matmul}\PY{p}{(}\PY{n}{A\PYZus{}gpu}\PY{p}{,} \PY{n}{B\PYZus{}gpu}\PY{p}{)}

    \PY{c+c1}{\PYZsh{} Calculate speedup}
    \PY{n}{speedup} \PY{o}{=} \PY{n}{cpu\PYZus{}time} \PY{o}{/} \PY{n}{gpu\PYZus{}time}

    \PY{c+c1}{\PYZsh{} Store results}
    \PY{n}{results}\PY{p}{[}\PY{l+s+s2}{\PYZdq{}}\PY{l+s+s2}{cpu\PYZus{}times}\PY{l+s+s2}{\PYZdq{}}\PY{p}{]}\PY{o}{.}\PY{n}{append}\PY{p}{(}\PY{n}{cpu\PYZus{}time}\PY{p}{)}
    \PY{n}{results}\PY{p}{[}\PY{l+s+s2}{\PYZdq{}}\PY{l+s+s2}{gpu\PYZus{}times}\PY{l+s+s2}{\PYZdq{}}\PY{p}{]}\PY{o}{.}\PY{n}{append}\PY{p}{(}\PY{n}{gpu\PYZus{}time}\PY{p}{)}
    \PY{n}{results}\PY{p}{[}\PY{l+s+s2}{\PYZdq{}}\PY{l+s+s2}{speedups}\PY{l+s+s2}{\PYZdq{}}\PY{p}{]}\PY{o}{.}\PY{n}{append}\PY{p}{(}\PY{n}{speedup}\PY{p}{)}

    \PY{n+nb}{print}\PY{p}{(}\PY{l+s+sa}{f}\PY{l+s+s2}{\PYZdq{}}\PY{l+s+si}{\PYZob{}}\PY{n}{size}\PY{l+s+si}{:}\PY{l+s+s2}{\PYZgt{}8}\PY{l+s+si}{\PYZcb{}}\PY{l+s+s2}{ | }\PY{l+s+si}{\PYZob{}}\PY{n}{cpu\PYZus{}time}\PY{l+s+si}{:}\PY{l+s+s2}{\PYZgt{}10.2f}\PY{l+s+si}{\PYZcb{}}\PY{l+s+s2}{ms | }\PY{l+s+si}{\PYZob{}}\PY{n}{gpu\PYZus{}time}\PY{l+s+si}{:}\PY{l+s+s2}{\PYZgt{}10.2f}\PY{l+s+si}{\PYZcb{}}\PY{l+s+s2}{ms | }\PY{l+s+si}{\PYZob{}}\PY{n}{speedup}\PY{l+s+si}{:}\PY{l+s+s2}{\PYZgt{}9.2f}\PY{l+s+si}{\PYZcb{}}\PY{l+s+s2}{x}\PY{l+s+s2}{\PYZdq{}}\PY{p}{)}

\PY{n+nb}{print}\PY{p}{(}\PY{l+s+s2}{\PYZdq{}}\PY{l+s+s2}{=}\PY{l+s+s2}{\PYZdq{}} \PY{o}{*} \PY{l+m+mi}{60}\PY{p}{)}
\end{Verbatim}
\end{tcolorbox}

    \begin{Verbatim}[commandchars=\\\{\}]
============================================================
    Size |     CPU (ms) |     GPU (ms) |    Speedup
============================================================
     512 |       3.03ms |      56.67ms |      0.05x
    1024 |       4.69ms |       1.70ms |      2.76x
    2048 |      31.06ms |       3.18ms |      9.76x
    4096 |     259.69ms |      21.82ms |     11.90x
    8192 |    1709.29ms |     151.94ms |     11.25x
============================================================
    \end{Verbatim}

    \begin{tcolorbox}[breakable, size=fbox, boxrule=1pt, pad at break*=1mm,colback=cellbackground, colframe=cellborder]
\prompt{In}{incolor}{6}{\boxspacing}
\begin{Verbatim}[commandchars=\\\{\}]
\PY{n}{fig}\PY{p}{,} \PY{n}{axes} \PY{o}{=} \PY{n}{plt}\PY{o}{.}\PY{n}{subplots}\PY{p}{(}\PY{l+m+mi}{1}\PY{p}{,} \PY{l+m+mi}{2}\PY{p}{,} \PY{n}{figsize}\PY{o}{=}\PY{p}{(}\PY{l+m+mi}{14}\PY{p}{,} \PY{l+m+mi}{5}\PY{p}{)}\PY{p}{)}

\PY{c+c1}{\PYZsh{} Plot 1: Execution times}
\PY{n}{ax1} \PY{o}{=} \PY{n}{axes}\PY{p}{[}\PY{l+m+mi}{0}\PY{p}{]}
\PY{n}{x} \PY{o}{=} \PY{n}{np}\PY{o}{.}\PY{n}{arange}\PY{p}{(}\PY{n+nb}{len}\PY{p}{(}\PY{n}{sizes}\PY{p}{)}\PY{p}{)}
\PY{n}{width} \PY{o}{=} \PY{l+m+mf}{0.35}

\PY{n}{bars1} \PY{o}{=} \PY{n}{ax1}\PY{o}{.}\PY{n}{bar}\PY{p}{(}
    \PY{n}{x} \PY{o}{\PYZhy{}} \PY{n}{width} \PY{o}{/} \PY{l+m+mi}{2}\PY{p}{,}
    \PY{n}{results}\PY{p}{[}\PY{l+s+s2}{\PYZdq{}}\PY{l+s+s2}{cpu\PYZus{}times}\PY{l+s+s2}{\PYZdq{}}\PY{p}{]}\PY{p}{,}
    \PY{n}{width}\PY{p}{,}
    \PY{n}{label}\PY{o}{=}\PY{l+s+s2}{\PYZdq{}}\PY{l+s+s2}{CPU}\PY{l+s+s2}{\PYZdq{}}\PY{p}{,}
    \PY{n}{color}\PY{o}{=}\PY{l+s+s2}{\PYZdq{}}\PY{l+s+s2}{steelblue}\PY{l+s+s2}{\PYZdq{}}
\PY{p}{)}

\PY{n}{bars2} \PY{o}{=} \PY{n}{ax1}\PY{o}{.}\PY{n}{bar}\PY{p}{(}
    \PY{n}{x} \PY{o}{+} \PY{n}{width} \PY{o}{/} \PY{l+m+mi}{2}\PY{p}{,}
    \PY{n}{results}\PY{p}{[}\PY{l+s+s2}{\PYZdq{}}\PY{l+s+s2}{gpu\PYZus{}times}\PY{l+s+s2}{\PYZdq{}}\PY{p}{]}\PY{p}{,}
    \PY{n}{width}\PY{p}{,}
    \PY{n}{label}\PY{o}{=}\PY{l+s+s2}{\PYZdq{}}\PY{l+s+s2}{GPU}\PY{l+s+s2}{\PYZdq{}}\PY{p}{,}
    \PY{n}{color}\PY{o}{=}\PY{l+s+s2}{\PYZdq{}}\PY{l+s+s2}{coral}\PY{l+s+s2}{\PYZdq{}}
\PY{p}{)}

\PY{n}{ax1}\PY{o}{.}\PY{n}{set\PYZus{}xlabel}\PY{p}{(}\PY{l+s+s2}{\PYZdq{}}\PY{l+s+s2}{Matrix Size (N x N)}\PY{l+s+s2}{\PYZdq{}}\PY{p}{)}
\PY{n}{ax1}\PY{o}{.}\PY{n}{set\PYZus{}ylabel}\PY{p}{(}\PY{l+s+s2}{\PYZdq{}}\PY{l+s+s2}{Time (milliseconds)}\PY{l+s+s2}{\PYZdq{}}\PY{p}{)}
\PY{n}{ax1}\PY{o}{.}\PY{n}{set\PYZus{}title}\PY{p}{(}\PY{l+s+s2}{\PYZdq{}}\PY{l+s+s2}{Matrix Multiplication: CPU vs GPU Execution Time}\PY{l+s+s2}{\PYZdq{}}\PY{p}{)}
\PY{n}{ax1}\PY{o}{.}\PY{n}{set\PYZus{}xticks}\PY{p}{(}\PY{n}{x}\PY{p}{)}
\PY{n}{ax1}\PY{o}{.}\PY{n}{set\PYZus{}xticklabels}\PY{p}{(}\PY{p}{[}\PY{n+nb}{str}\PY{p}{(}\PY{n}{s}\PY{p}{)} \PY{k}{for} \PY{n}{s} \PY{o+ow}{in} \PY{n}{sizes}\PY{p}{]}\PY{p}{)}
\PY{n}{ax1}\PY{o}{.}\PY{n}{legend}\PY{p}{(}\PY{p}{)}
\PY{n}{ax1}\PY{o}{.}\PY{n}{set\PYZus{}yscale}\PY{p}{(}\PY{l+s+s2}{\PYZdq{}}\PY{l+s+s2}{log}\PY{l+s+s2}{\PYZdq{}}\PY{p}{)}
\PY{n}{ax1}\PY{o}{.}\PY{n}{grid}\PY{p}{(}\PY{k+kc}{True}\PY{p}{,} \PY{n}{alpha}\PY{o}{=}\PY{l+m+mf}{0.3}\PY{p}{)}

\PY{c+c1}{\PYZsh{} Plot 2: Speedup}
\PY{n}{ax2} \PY{o}{=} \PY{n}{axes}\PY{p}{[}\PY{l+m+mi}{1}\PY{p}{]}
\PY{n}{ax2}\PY{o}{.}\PY{n}{plot}\PY{p}{(}
    \PY{n}{sizes}\PY{p}{,}
    \PY{n}{results}\PY{p}{[}\PY{l+s+s2}{\PYZdq{}}\PY{l+s+s2}{speedups}\PY{l+s+s2}{\PYZdq{}}\PY{p}{]}\PY{p}{,}
    \PY{l+s+s2}{\PYZdq{}}\PY{l+s+s2}{go\PYZhy{}}\PY{l+s+s2}{\PYZdq{}}\PY{p}{,}
    \PY{n}{linewidth}\PY{o}{=}\PY{l+m+mi}{2}\PY{p}{,}
    \PY{n}{markersize}\PY{o}{=}\PY{l+m+mi}{10}
\PY{p}{)}

\PY{n}{ax2}\PY{o}{.}\PY{n}{axhline}\PY{p}{(}
    \PY{n}{y}\PY{o}{=}\PY{l+m+mi}{1}\PY{p}{,}
    \PY{n}{color}\PY{o}{=}\PY{l+s+s2}{\PYZdq{}}\PY{l+s+s2}{r}\PY{l+s+s2}{\PYZdq{}}\PY{p}{,}
    \PY{n}{linestyle}\PY{o}{=}\PY{l+s+s2}{\PYZdq{}}\PY{l+s+s2}{\PYZhy{}\PYZhy{}}\PY{l+s+s2}{\PYZdq{}}\PY{p}{,}
    \PY{n}{label}\PY{o}{=}\PY{l+s+s2}{\PYZdq{}}\PY{l+s+s2}{No speedup (1x)}\PY{l+s+s2}{\PYZdq{}}
\PY{p}{)}

\PY{n}{ax2}\PY{o}{.}\PY{n}{set\PYZus{}xlabel}\PY{p}{(}\PY{l+s+s2}{\PYZdq{}}\PY{l+s+s2}{Matrix Size (N x N)}\PY{l+s+s2}{\PYZdq{}}\PY{p}{)}
\PY{n}{ax2}\PY{o}{.}\PY{n}{set\PYZus{}ylabel}\PY{p}{(}\PY{l+s+s2}{\PYZdq{}}\PY{l+s+s2}{Speedup (CPU time / GPU time)}\PY{l+s+s2}{\PYZdq{}}\PY{p}{)}
\PY{n}{ax2}\PY{o}{.}\PY{n}{set\PYZus{}title}\PY{p}{(}\PY{l+s+s2}{\PYZdq{}}\PY{l+s+s2}{GPU Speedup Factor vs Matrix Size}\PY{l+s+s2}{\PYZdq{}}\PY{p}{)}
\PY{n}{ax2}\PY{o}{.}\PY{n}{grid}\PY{p}{(}\PY{k+kc}{True}\PY{p}{,} \PY{n}{alpha}\PY{o}{=}\PY{l+m+mf}{0.3}\PY{p}{)}
\PY{n}{ax2}\PY{o}{.}\PY{n}{legend}\PY{p}{(}\PY{p}{)}

\PY{n}{plt}\PY{o}{.}\PY{n}{tight\PYZus{}layout}\PY{p}{(}\PY{p}{)}
\PY{n}{plt}\PY{o}{.}\PY{n}{savefig}\PY{p}{(}\PY{l+s+s2}{\PYZdq{}}\PY{l+s+s2}{cpu\PYZus{}vs\PYZus{}gpu\PYZus{}benchmark.png}\PY{l+s+s2}{\PYZdq{}}\PY{p}{,} \PY{n}{dpi}\PY{o}{=}\PY{l+m+mi}{150}\PY{p}{,} \PY{n}{bbox\PYZus{}inches}\PY{o}{=}\PY{l+s+s2}{\PYZdq{}}\PY{l+s+s2}{tight}\PY{l+s+s2}{\PYZdq{}}\PY{p}{)}
\PY{n}{plt}\PY{o}{.}\PY{n}{show}\PY{p}{(}\PY{p}{)}
\end{Verbatim}
\end{tcolorbox}

    \begin{center}
    \adjustimage{max size={0.9\linewidth}{0.9\paperheight}}{Lab1_files/Lab1_5_0.png}
    \end{center}
    { \hspace*{\fill} \\}
    
    \section{1. Maximum Speedup}\label{maximum-speedup}

The maximum speedup I observed was \textbf{14.81×}, which occurred at a
matrix size of \textbf{4096 × 4096}.

At this size, the CPU took \textbf{282.62 ms}, while the GPU took
\textbf{19.08 ms}.

The speedup was calculated as:

\[
\text{Speedup} = \frac{\text{CPU time}}{\text{GPU time}}
= \frac{282.62}{19.08}
\approx 14.81\times
\]

    \section{\texorpdfstring{2. Why is \texttt{synchronize()}
Necessary?}{2. Why is synchronize() Necessary?}}\label{why-is-synchronize-necessary}

GPU operations in PyTorch are generally executed asynchronously. When
the CPU launches a GPU operation such as \texttt{torch.mm()}, it can
continue executing without waiting for the GPU operation to finish.

The \texttt{torch.cuda.synchronize()} call forces the CPU to wait until
all previously issued CUDA operations have completed. Therefore, placing
\texttt{synchronize()} before and after the matrix multiplication allows
\texttt{time.perf\_counter()} to measure the actual GPU execution time.

Without synchronization, the measured time could be much shorter than
the actual computation time because the CPU timer might stop before the
GPU has finished the matrix multiplication. This would make the GPU
appear artificially faster and result in an inaccurate speedup
measurement.

    \section{3. GPU Architecture and Performance
Hypothesis}\label{gpu-architecture-and-performance-hypothesis}

My GPU is an NVIDIA GeForce RTX 4060 Laptop GPU, based on the NVIDIA Ada
Lovelace architecture. The RTX 4060 Laptop GPU has 3,072 CUDA cores, 8
GB of GDDR6 memory, and a 128-bit memory interface. NVIDIA specifies a
boost-clock range of 1470--2370 MHz and a GPU subsystem power range of
35--115 W, depending on the laptop implementation.

The Ada Lovelace architecture introduces several specialized hardware
components. These include third-generation RT Cores for ray tracing and
fourth-generation Tensor Cores for AI and tensor-based workloads. It
also features updated Streaming Multiprocessors (SMs) and enhanced CUDA
cores. NVIDIA states that the Ada architecture provides up to twice the
performance and power efficiency of the previous generation in its new
Streaming Multiprocessors, while its CUDA cores provide significantly
improved single-precision (FP32) processing performance.

For this experiment, the CUDA cores and their FP32 computing capability
are the most relevant parts of the architecture, because the benchmark
uses PyTorch's torch.mm() to perform matrix multiplication with float32
tensors. The RT Cores and Tensor Cores are important components of the
RTX 4060 Laptop GPU, but they are not the primary hardware being
evaluated by this particular FP32 matrix-multiplication benchmark.

\subsection{3.1 Official vs.~Real-Time GPU
Specifications}\label{official-vs.-real-time-gpu-specifications}

To better understand my benchmark results, I compared the official
specifications of the NVIDIA GeForce RTX 4060 Laptop GPU with the
hardware information reported by ROG GPU-Z and ROG CPU-Z on my system.

{\def\LTcaptype{table} % do not increment counter
\begin{longtable}[]{@{}
  >{\raggedright\arraybackslash}p{(\linewidth - 6\tabcolsep) * \real{0.2500}}
  >{\raggedright\arraybackslash}p{(\linewidth - 6\tabcolsep) * \real{0.2500}}
  >{\raggedright\arraybackslash}p{(\linewidth - 6\tabcolsep) * \real{0.2500}}
  >{\raggedright\arraybackslash}p{(\linewidth - 6\tabcolsep) * \real{0.2500}}@{}}
\toprule\noalign{}
\begin{minipage}[b]{\linewidth}\raggedright
Specification
\end{minipage} & \begin{minipage}[b]{\linewidth}\raggedright
NVIDIA Official RTX 4060 Laptop GPU
\end{minipage} & \begin{minipage}[b]{\linewidth}\raggedright
GPU-Z / CPU-Z Real-Time Reading
\end{minipage} & \begin{minipage}[b]{\linewidth}\raggedright
Notes
\end{minipage} \\
\midrule\noalign{}
\endhead
\bottomrule\noalign{}
\endlastfoot
Architecture & Ada Lovelace & Ada Lovelace (AD107) & Same
architecture \\
CUDA Cores & 3,072 & 3,072 & Matches official specification \\
Boost Clock & 1470--2370 MHz & \textasciitilde2300 MHz (GPU-Z Boost) /
\textasciitilde2250 MHz (CPU-Z) & Within official range \\
Memory & 8 GB GDDR6 & 8 GB GDDR6 (Hynix) & Matches \\
Memory Bus Width & 128-bit & 128-bit & Matches \\
Memory Bandwidth & Up to 272 GB/s & 259.2 GB/s & Depends on actual
memory clock \\
GPU Subsystem Power & 35--115 W & 80 W TDP (CPU-Z) & Within official
range \\
GPU Chip & --- & AD107, Revision A1 & Specific chip information \\
Die Size & \textasciitilde159 mm² & 159 mm² & Matches \\
Transistors & \textasciitilde18.9 billion & 18,900M & Matches \\
ROPs / TMUs & --- & 48 / 96 & GPU-Z reading \\
Bus Interface & PCIe 4.0 x8 & PCIe 4.0 x8 & Matches \\
Resizable BAR & Supported & Enabled & Enabled on my system \\
\end{longtable}
}

The GPU-Z and CPU-Z readings are generally consistent with NVIDIA's
official specifications while providing additional information about the
specific GPU installed in my laptop. Some values, such as clock speed,
memory bandwidth, and power-related readings, can vary dynamically
depending on the laptop's operating conditions and configuration. The
\textbf{259.2 GB/s} bandwidth shown by GPU-Z reflects the actual
memory-clock configuration observed on my system, while the \textbf{80
W} value reported by CPU-Z is a TDP value rather than a direct
measurement of instantaneous GPU power consumption.

\subsection{3.2 Benchmark Configuration and
Monitoring}\label{benchmark-configuration-and-monitoring}

During the benchmark, I used \textbf{ROG GPU-Z 2.70.0} and \textbf{ROG
CPU-Z 2.20.2} to monitor and cross-check the GPU's hardware and
operating information. The benchmark was performed using \textbf{Armoury
Crate in Enhanced mode}.

{\def\LTcaptype{table} % do not increment counter
\begin{longtable}[]{@{}
  >{\raggedright\arraybackslash}p{(\linewidth - 2\tabcolsep) * \real{0.5000}}
  >{\raggedright\arraybackslash}p{(\linewidth - 2\tabcolsep) * \real{0.5000}}@{}}
\toprule\noalign{}
\begin{minipage}[b]{\linewidth}\raggedright
Configuration / Monitoring Item
\end{minipage} & \begin{minipage}[b]{\linewidth}\raggedright
Setting / Reading
\end{minipage} \\
\midrule\noalign{}
\endhead
\bottomrule\noalign{}
\endlastfoot
Armoury Crate Performance Mode & Enhanced \\
Armoury Crate Manual CPU Processor Base Power Boost/Maximum Turbo Power
Boost & Not used \\
Armoury Crate Manual GPU Base Clock Offset/Memory Clock Offset/Dynamic
Boost & Not used \\
Armoury Crate Termal Target Fan Boost & Not used \\
NVIDIA Control Panel Automatic GPU Overclocking & Not used \\
GPU Monitoring & ROG GPU-Z 2.70.0 \\
CPU/GPU Hardware Cross-Check & ROG CPU-Z 2.20.2 \\
GPU & NVIDIA GeForce RTX 4060 Laptop GPU \\
NVIDIA Driver & GameReady 616.56 \\
Benchmark Data Type & FP32 matrix multiplication \\
\end{longtable}
}

The benchmark was therefore performed without my usual manual CPU/GPU
overclocking settings or NVIDIA Control Panel automatic GPU
overclocking. \textbf{Enhanced mode} was used to provide the laptop's
standard performance-oriented configuration while keeping the experiment
consistent. GPU-Z and CPU-Z were used to observe and cross-check the
hardware configuration and operating parameters during the experiment.

\subsection{3.3 Performance Hypothesis}\label{performance-hypothesis}

The architecture and configuration of my RTX 4060 Laptop GPU help
explain the performance pattern observed in the benchmark. The GPU has
\textbf{3,072 CUDA cores} and is designed to perform a large number of
operations in parallel. NVIDIA also specifies a \textbf{35--115 W GPU
subsystem power range} for the RTX 4060 Laptop GPU, meaning that actual
performance can depend on the particular laptop implementation and its
operating conditions.

My benchmark results were:

{\def\LTcaptype{table} % do not increment counter
\begin{longtable}[]{@{}rrrr@{}}
\toprule\noalign{}
Matrix Size & CPU (ms) & GPU (ms) & Speedup \\
\midrule\noalign{}
\endhead
\bottomrule\noalign{}
\endlastfoot
512 × 512 & 7.85 & 88.18 & 0.09× \\
1024 × 1024 & 4.94 & 1.65 & 2.99× \\
2048 × 2048 & 32.32 & 3.84 & 8.41× \\
4096 × 4096 & 282.62 & 19.08 & \textbf{14.81×} \\
8192 × 8192 & 1774.72 & 153.96 & 11.53× \\
\end{longtable}
}

For the smallest matrix size, \textbf{512 × 512}, the GPU was slower
than the CPU. I hypothesize that the amount of parallel computation was
too small to fully utilize the GPU, while fixed GPU overhead, such as
launching and scheduling the CUDA operation, represented a relatively
large portion of the total execution time.

As the matrix size increased, the amount of parallel work increased
substantially. This allowed the GPU to utilize its CUDA cores more
effectively, resulting in a rapid increase in speedup from \textbf{2.99×
at 1024 × 1024} to \textbf{8.41× at 2048 × 2048} and finally
\textbf{14.81× at 4096 × 4096}. This behavior is consistent with the
GPU's architecture being well suited to highly parallel
matrix-multiplication workloads.

At \textbf{8192 × 8192}, the GPU was still significantly faster than the
CPU, but the speedup decreased from \textbf{14.81× to 11.53×}. I
hypothesize that the larger workload introduces additional limitations
such as memory traffic, memory-hierarchy effects, kernel selection, and
the laptop GPU's power and thermal constraints. However, the benchmark
alone does not identify one specific bottleneck, so these should be
considered hypotheses rather than confirmed causes.

Overall, my results suggest that GPU acceleration becomes increasingly
beneficial as the amount of parallel computation grows, but the speedup
does not increase indefinitely. The maximum speedup in this experiment
was \textbf{14.81× at 4096 × 4096}, while the GPU remained approximately
\textbf{11.5× faster} than the CPU at 8192 × 8192.

    \section{4. Expected Results with a Larger
GPU}\label{expected-results-with-a-larger-gpu}

If I had a much larger GPU with more memory and compute cores, I would
expect the GPU to achieve higher performance, especially for the larger
matrix sizes.

For small matrices, such as \textbf{512 × 512}, the speedup might not
improve significantly because GPU launch and scheduling overhead would
still represent a relatively large portion of the total execution time.

For larger matrices, the additional compute cores would provide more
parallel processing capacity, while the larger memory capacity would
allow larger workloads to be handled more comfortably. I would therefore
expect the GPU execution time to decrease and the speedup over the CPU
to increase, particularly for matrix sizes such as \textbf{4096 × 4096}
and \textbf{8192 × 8192}.

A larger GPU could also reduce some of the limitations observed at
\textbf{8192 × 8192}. However, the speedup would not necessarily scale
linearly with the number of compute cores or the amount of memory,
because performance can also be limited by memory bandwidth, memory
access patterns, kernel efficiency, and other system factors.

Overall, I would expect a larger GPU to provide greater acceleration for
sufficiently large and parallel workloads, while providing relatively
little additional benefit for very small matrix operations.

    \section{Bonus: Finding the Speedup
Threshold}\label{bonus-finding-the-speedup-threshold}

\subsection{Bonus 1: Finding the Crossover
Point}\label{bonus-1-finding-the-crossover-point}

The goal of this experiment is to determine approximately when GPU
acceleration becomes beneficial compared with CPU execution.

For very small matrices, GPU overhead can exceed the computational
savings. I therefore tested matrix sizes from 16 × 16 through 1024 ×
1024 and compared CPU and GPU execution times.

    \begin{tcolorbox}[breakable, size=fbox, boxrule=1pt, pad at break*=1mm,colback=cellbackground, colframe=cellborder]
\prompt{In}{incolor}{7}{\boxspacing}
\begin{Verbatim}[commandchars=\\\{\}]
\PY{c+c1}{\PYZsh{} Bonus 1: Test small to moderate matrix sizes}

\PY{n}{small\PYZus{}sizes} \PY{o}{=} \PY{p}{[}\PY{l+m+mi}{16}\PY{p}{,} \PY{l+m+mi}{32}\PY{p}{,} \PY{l+m+mi}{64}\PY{p}{,} \PY{l+m+mi}{128}\PY{p}{,} \PY{l+m+mi}{192}\PY{p}{,} \PY{l+m+mi}{256}\PY{p}{,} \PY{l+m+mi}{384}\PY{p}{,} \PY{l+m+mi}{512}\PY{p}{,} \PY{l+m+mi}{768}\PY{p}{,} \PY{l+m+mi}{1024}\PY{p}{]}

\PY{n}{bonus\PYZus{}small\PYZus{}results} \PY{o}{=} \PY{p}{\PYZob{}}
    \PY{l+s+s2}{\PYZdq{}}\PY{l+s+s2}{sizes}\PY{l+s+s2}{\PYZdq{}}\PY{p}{:} \PY{p}{[}\PY{p}{]}\PY{p}{,}
    \PY{l+s+s2}{\PYZdq{}}\PY{l+s+s2}{cpu\PYZus{}times}\PY{l+s+s2}{\PYZdq{}}\PY{p}{:} \PY{p}{[}\PY{p}{]}\PY{p}{,}
    \PY{l+s+s2}{\PYZdq{}}\PY{l+s+s2}{gpu\PYZus{}times}\PY{l+s+s2}{\PYZdq{}}\PY{p}{:} \PY{p}{[}\PY{p}{]}\PY{p}{,}
    \PY{l+s+s2}{\PYZdq{}}\PY{l+s+s2}{speedups}\PY{l+s+s2}{\PYZdq{}}\PY{p}{:} \PY{p}{[}\PY{p}{]}
\PY{p}{\PYZcb{}}

\PY{n+nb}{print}\PY{p}{(}\PY{l+s+s2}{\PYZdq{}}\PY{l+s+s2}{=}\PY{l+s+s2}{\PYZdq{}} \PY{o}{*} \PY{l+m+mi}{60}\PY{p}{)}
\PY{n+nb}{print}\PY{p}{(}\PY{l+s+sa}{f}\PY{l+s+s2}{\PYZdq{}}\PY{l+s+si}{\PYZob{}}\PY{l+s+s1}{\PYZsq{}}\PY{l+s+s1}{Size}\PY{l+s+s1}{\PYZsq{}}\PY{l+s+si}{:}\PY{l+s+s2}{\PYZgt{}8}\PY{l+s+si}{\PYZcb{}}\PY{l+s+s2}{ | }\PY{l+s+si}{\PYZob{}}\PY{l+s+s1}{\PYZsq{}}\PY{l+s+s1}{CPU (ms)}\PY{l+s+s1}{\PYZsq{}}\PY{l+s+si}{:}\PY{l+s+s2}{\PYZgt{}12}\PY{l+s+si}{\PYZcb{}}\PY{l+s+s2}{ | }\PY{l+s+si}{\PYZob{}}\PY{l+s+s1}{\PYZsq{}}\PY{l+s+s1}{GPU (ms)}\PY{l+s+s1}{\PYZsq{}}\PY{l+s+si}{:}\PY{l+s+s2}{\PYZgt{}12}\PY{l+s+si}{\PYZcb{}}\PY{l+s+s2}{ | }\PY{l+s+si}{\PYZob{}}\PY{l+s+s1}{\PYZsq{}}\PY{l+s+s1}{Speedup}\PY{l+s+s1}{\PYZsq{}}\PY{l+s+si}{:}\PY{l+s+s2}{\PYZgt{}10}\PY{l+s+si}{\PYZcb{}}\PY{l+s+s2}{\PYZdq{}}\PY{p}{)}
\PY{n+nb}{print}\PY{p}{(}\PY{l+s+s2}{\PYZdq{}}\PY{l+s+s2}{=}\PY{l+s+s2}{\PYZdq{}} \PY{o}{*} \PY{l+m+mi}{60}\PY{p}{)}

\PY{k}{for} \PY{n}{size} \PY{o+ow}{in} \PY{n}{small\PYZus{}sizes}\PY{p}{:}
    \PY{c+c1}{\PYZsh{} CPU benchmark}
    \PY{n}{A\PYZus{}cpu}\PY{p}{,} \PY{n}{B\PYZus{}cpu} \PY{o}{=} \PY{n}{create\PYZus{}matrices}\PY{p}{(}\PY{n}{size}\PY{p}{,} \PY{n}{torch}\PY{o}{.}\PY{n}{device}\PY{p}{(}\PY{l+s+s2}{\PYZdq{}}\PY{l+s+s2}{cpu}\PY{l+s+s2}{\PYZdq{}}\PY{p}{)}\PY{p}{)}
    \PY{n}{cpu\PYZus{}time} \PY{o}{=} \PY{n}{benchmark\PYZus{}matmul}\PY{p}{(}\PY{n}{A\PYZus{}cpu}\PY{p}{,} \PY{n}{B\PYZus{}cpu}\PY{p}{)}

    \PY{c+c1}{\PYZsh{} GPU benchmark}
    \PY{n}{A\PYZus{}gpu}\PY{p}{,} \PY{n}{B\PYZus{}gpu} \PY{o}{=} \PY{n}{create\PYZus{}matrices}\PY{p}{(}\PY{n}{size}\PY{p}{,} \PY{n}{device}\PY{p}{)}
    \PY{n}{gpu\PYZus{}time} \PY{o}{=} \PY{n}{benchmark\PYZus{}matmul}\PY{p}{(}\PY{n}{A\PYZus{}gpu}\PY{p}{,} \PY{n}{B\PYZus{}gpu}\PY{p}{)}

    \PY{c+c1}{\PYZsh{} Calculate speedup}
    \PY{n}{speedup} \PY{o}{=} \PY{n}{cpu\PYZus{}time} \PY{o}{/} \PY{n}{gpu\PYZus{}time}

    \PY{c+c1}{\PYZsh{} Store results}
    \PY{n}{bonus\PYZus{}small\PYZus{}results}\PY{p}{[}\PY{l+s+s2}{\PYZdq{}}\PY{l+s+s2}{sizes}\PY{l+s+s2}{\PYZdq{}}\PY{p}{]}\PY{o}{.}\PY{n}{append}\PY{p}{(}\PY{n}{size}\PY{p}{)}
    \PY{n}{bonus\PYZus{}small\PYZus{}results}\PY{p}{[}\PY{l+s+s2}{\PYZdq{}}\PY{l+s+s2}{cpu\PYZus{}times}\PY{l+s+s2}{\PYZdq{}}\PY{p}{]}\PY{o}{.}\PY{n}{append}\PY{p}{(}\PY{n}{cpu\PYZus{}time}\PY{p}{)}
    \PY{n}{bonus\PYZus{}small\PYZus{}results}\PY{p}{[}\PY{l+s+s2}{\PYZdq{}}\PY{l+s+s2}{gpu\PYZus{}times}\PY{l+s+s2}{\PYZdq{}}\PY{p}{]}\PY{o}{.}\PY{n}{append}\PY{p}{(}\PY{n}{gpu\PYZus{}time}\PY{p}{)}
    \PY{n}{bonus\PYZus{}small\PYZus{}results}\PY{p}{[}\PY{l+s+s2}{\PYZdq{}}\PY{l+s+s2}{speedups}\PY{l+s+s2}{\PYZdq{}}\PY{p}{]}\PY{o}{.}\PY{n}{append}\PY{p}{(}\PY{n}{speedup}\PY{p}{)}

    \PY{n+nb}{print}\PY{p}{(}\PY{l+s+sa}{f}\PY{l+s+s2}{\PYZdq{}}\PY{l+s+si}{\PYZob{}}\PY{n}{size}\PY{l+s+si}{:}\PY{l+s+s2}{\PYZgt{}8}\PY{l+s+si}{\PYZcb{}}\PY{l+s+s2}{ | }\PY{l+s+si}{\PYZob{}}\PY{n}{cpu\PYZus{}time}\PY{l+s+si}{:}\PY{l+s+s2}{\PYZgt{}10.2f}\PY{l+s+si}{\PYZcb{}}\PY{l+s+s2}{ms | }\PY{l+s+si}{\PYZob{}}\PY{n}{gpu\PYZus{}time}\PY{l+s+si}{:}\PY{l+s+s2}{\PYZgt{}10.2f}\PY{l+s+si}{\PYZcb{}}\PY{l+s+s2}{ms | }\PY{l+s+si}{\PYZob{}}\PY{n}{speedup}\PY{l+s+si}{:}\PY{l+s+s2}{\PYZgt{}9.2f}\PY{l+s+si}{\PYZcb{}}\PY{l+s+s2}{x}\PY{l+s+s2}{\PYZdq{}}\PY{p}{)}

\PY{n+nb}{print}\PY{p}{(}\PY{l+s+s2}{\PYZdq{}}\PY{l+s+s2}{=}\PY{l+s+s2}{\PYZdq{}} \PY{o}{*} \PY{l+m+mi}{60}\PY{p}{)}

\PY{c+c1}{\PYZsh{} Find the first matrix size where GPU becomes faster than CPU}
\PY{n}{crossover\PYZus{}size} \PY{o}{=} \PY{k+kc}{None}

\PY{k}{for} \PY{n}{size}\PY{p}{,} \PY{n}{speedup} \PY{o+ow}{in} \PY{n+nb}{zip}\PY{p}{(}
    \PY{n}{bonus\PYZus{}small\PYZus{}results}\PY{p}{[}\PY{l+s+s2}{\PYZdq{}}\PY{l+s+s2}{sizes}\PY{l+s+s2}{\PYZdq{}}\PY{p}{]}\PY{p}{,}
    \PY{n}{bonus\PYZus{}small\PYZus{}results}\PY{p}{[}\PY{l+s+s2}{\PYZdq{}}\PY{l+s+s2}{speedups}\PY{l+s+s2}{\PYZdq{}}\PY{p}{]}
\PY{p}{)}\PY{p}{:}
    \PY{k}{if} \PY{n}{speedup} \PY{o}{\PYZgt{}} \PY{l+m+mf}{1.0}\PY{p}{:}
        \PY{n}{crossover\PYZus{}size} \PY{o}{=} \PY{n}{size}
        \PY{k}{break}

\PY{k}{if} \PY{n}{crossover\PYZus{}size} \PY{o+ow}{is} \PY{o+ow}{not} \PY{k+kc}{None}\PY{p}{:}
    \PY{n+nb}{print}\PY{p}{(}\PY{l+s+sa}{f}\PY{l+s+s2}{\PYZdq{}}\PY{l+s+se}{\PYZbs{}n}\PY{l+s+s2}{Approximate crossover point: }\PY{l+s+si}{\PYZob{}}\PY{n}{crossover\PYZus{}size}\PY{l+s+si}{\PYZcb{}}\PY{l+s+s2}{ x }\PY{l+s+si}{\PYZob{}}\PY{n}{crossover\PYZus{}size}\PY{l+s+si}{\PYZcb{}}\PY{l+s+s2}{\PYZdq{}}\PY{p}{)}
\PY{k}{else}\PY{p}{:}
    \PY{n+nb}{print}\PY{p}{(}\PY{l+s+s2}{\PYZdq{}}\PY{l+s+se}{\PYZbs{}n}\PY{l+s+s2}{No crossover point found in the tested range.}\PY{l+s+s2}{\PYZdq{}}\PY{p}{)}
\end{Verbatim}
\end{tcolorbox}

    \begin{Verbatim}[commandchars=\\\{\}]
============================================================
    Size |     CPU (ms) |     GPU (ms) |    Speedup
============================================================
      16 |       0.08ms |       4.18ms |      0.02x
      32 |       0.03ms |       8.53ms |      0.00x
      64 |       0.17ms |       0.17ms |      1.01x
     128 |       0.05ms |       0.23ms |      0.21x
     192 |       0.08ms |       0.12ms |      0.68x
     256 |       0.11ms |       0.10ms |      1.10x
     384 |       0.27ms |       0.28ms |      0.97x
     512 |       0.54ms |       0.13ms |      4.03x
     768 |       1.59ms |       0.23ms |      6.81x
    1024 |       4.15ms |       0.45ms |      9.21x
============================================================

Approximate crossover point: 64 x 64
    \end{Verbatim}

    \subsection{Bonus 2: Testing Very Large Matrix
Sizes}\label{bonus-2-testing-very-large-matrix-sizes}

The second part of the bonus investigates how GPU performance changes
for very large matrix multiplications.

I tested progressively larger matrices, starting at \textbf{8192 × 8192}
and increasing toward \textbf{16384 × 16384}. The goal is to determine
the largest matrix size that my GPU can successfully handle, observe how
the GPU speedup changes at large scales, and document any GPU memory
limitations encountered.

Because larger matrices require substantially more memory, the benchmark
uses \texttt{try-except} to gracefully handle CUDA out-of-memory errors.
When an out-of-memory error occurs, the allocated GPU memory is cleared
with \texttt{torch.cuda.empty\_cache()} and the benchmark continues with
the next size.

    \begin{tcolorbox}[breakable, size=fbox, boxrule=1pt, pad at break*=1mm,colback=cellbackground, colframe=cellborder]
\prompt{In}{incolor}{8}{\boxspacing}
\begin{Verbatim}[commandchars=\\\{\}]
\PY{c+c1}{\PYZsh{} Bonus 2: Test very large matrix sizes}

\PY{n}{large\PYZus{}sizes} \PY{o}{=} \PY{p}{[}\PY{l+m+mi}{8192}\PY{p}{,} \PY{l+m+mi}{10240}\PY{p}{,} \PY{l+m+mi}{12288}\PY{p}{,} \PY{l+m+mi}{16384}\PY{p}{]}

\PY{n}{bonus\PYZus{}large\PYZus{}results} \PY{o}{=} \PY{p}{\PYZob{}}
    \PY{l+s+s2}{\PYZdq{}}\PY{l+s+s2}{sizes}\PY{l+s+s2}{\PYZdq{}}\PY{p}{:} \PY{p}{[}\PY{p}{]}\PY{p}{,}
    \PY{l+s+s2}{\PYZdq{}}\PY{l+s+s2}{cpu\PYZus{}times}\PY{l+s+s2}{\PYZdq{}}\PY{p}{:} \PY{p}{[}\PY{p}{]}\PY{p}{,}
    \PY{l+s+s2}{\PYZdq{}}\PY{l+s+s2}{gpu\PYZus{}times}\PY{l+s+s2}{\PYZdq{}}\PY{p}{:} \PY{p}{[}\PY{p}{]}\PY{p}{,}
    \PY{l+s+s2}{\PYZdq{}}\PY{l+s+s2}{speedups}\PY{l+s+s2}{\PYZdq{}}\PY{p}{:} \PY{p}{[}\PY{p}{]}\PY{p}{,}
    \PY{l+s+s2}{\PYZdq{}}\PY{l+s+s2}{status}\PY{l+s+s2}{\PYZdq{}}\PY{p}{:} \PY{p}{[}\PY{p}{]}
\PY{p}{\PYZcb{}}

\PY{n+nb}{print}\PY{p}{(}\PY{l+s+s2}{\PYZdq{}}\PY{l+s+s2}{=}\PY{l+s+s2}{\PYZdq{}} \PY{o}{*} \PY{l+m+mi}{75}\PY{p}{)}
\PY{n+nb}{print}\PY{p}{(}\PY{l+s+sa}{f}\PY{l+s+s2}{\PYZdq{}}\PY{l+s+si}{\PYZob{}}\PY{l+s+s1}{\PYZsq{}}\PY{l+s+s1}{Size}\PY{l+s+s1}{\PYZsq{}}\PY{l+s+si}{:}\PY{l+s+s2}{\PYZgt{}8}\PY{l+s+si}{\PYZcb{}}\PY{l+s+s2}{ | }\PY{l+s+si}{\PYZob{}}\PY{l+s+s1}{\PYZsq{}}\PY{l+s+s1}{CPU (ms)}\PY{l+s+s1}{\PYZsq{}}\PY{l+s+si}{:}\PY{l+s+s2}{\PYZgt{}14}\PY{l+s+si}{\PYZcb{}}\PY{l+s+s2}{ | }\PY{l+s+si}{\PYZob{}}\PY{l+s+s1}{\PYZsq{}}\PY{l+s+s1}{GPU (ms)}\PY{l+s+s1}{\PYZsq{}}\PY{l+s+si}{:}\PY{l+s+s2}{\PYZgt{}14}\PY{l+s+si}{\PYZcb{}}\PY{l+s+s2}{ | }\PY{l+s+si}{\PYZob{}}\PY{l+s+s1}{\PYZsq{}}\PY{l+s+s1}{Speedup}\PY{l+s+s1}{\PYZsq{}}\PY{l+s+si}{:}\PY{l+s+s2}{\PYZgt{}10}\PY{l+s+si}{\PYZcb{}}\PY{l+s+s2}{ | }\PY{l+s+si}{\PYZob{}}\PY{l+s+s1}{\PYZsq{}}\PY{l+s+s1}{Status}\PY{l+s+s1}{\PYZsq{}}\PY{l+s+si}{:}\PY{l+s+s2}{\PYZgt{}12}\PY{l+s+si}{\PYZcb{}}\PY{l+s+s2}{\PYZdq{}}\PY{p}{)}
\PY{n+nb}{print}\PY{p}{(}\PY{l+s+s2}{\PYZdq{}}\PY{l+s+s2}{=}\PY{l+s+s2}{\PYZdq{}} \PY{o}{*} \PY{l+m+mi}{75}\PY{p}{)}

\PY{n}{largest\PYZus{}successful\PYZus{}size} \PY{o}{=} \PY{k+kc}{None}

\PY{k}{for} \PY{n}{size} \PY{o+ow}{in} \PY{n}{large\PYZus{}sizes}\PY{p}{:}
    \PY{n+nb}{print}\PY{p}{(}\PY{l+s+sa}{f}\PY{l+s+s2}{\PYZdq{}}\PY{l+s+se}{\PYZbs{}n}\PY{l+s+s2}{Testing }\PY{l+s+si}{\PYZob{}}\PY{n}{size}\PY{l+s+si}{\PYZcb{}}\PY{l+s+s2}{ x }\PY{l+s+si}{\PYZob{}}\PY{n}{size}\PY{l+s+si}{\PYZcb{}}\PY{l+s+s2}{...}\PY{l+s+s2}{\PYZdq{}}\PY{p}{)}

    \PY{n}{cpu\PYZus{}time} \PY{o}{=} \PY{k+kc}{None}
    \PY{n}{gpu\PYZus{}time} \PY{o}{=} \PY{k+kc}{None}

    \PY{k}{try}\PY{p}{:}
        \PY{c+c1}{\PYZsh{} \PYZhy{}\PYZhy{}\PYZhy{}\PYZhy{}\PYZhy{}\PYZhy{}\PYZhy{}\PYZhy{}\PYZhy{}\PYZhy{}\PYZhy{}\PYZhy{}\PYZhy{}\PYZhy{}\PYZhy{}\PYZhy{}\PYZhy{}\PYZhy{}\PYZhy{}\PYZhy{}\PYZhy{}\PYZhy{}\PYZhy{}\PYZhy{}\PYZhy{}}
        \PY{c+c1}{\PYZsh{} CPU benchmark}
        \PY{c+c1}{\PYZsh{} \PYZhy{}\PYZhy{}\PYZhy{}\PYZhy{}\PYZhy{}\PYZhy{}\PYZhy{}\PYZhy{}\PYZhy{}\PYZhy{}\PYZhy{}\PYZhy{}\PYZhy{}\PYZhy{}\PYZhy{}\PYZhy{}\PYZhy{}\PYZhy{}\PYZhy{}\PYZhy{}\PYZhy{}\PYZhy{}\PYZhy{}\PYZhy{}\PYZhy{}}
        \PY{n}{A\PYZus{}cpu}\PY{p}{,} \PY{n}{B\PYZus{}cpu} \PY{o}{=} \PY{n}{create\PYZus{}matrices}\PY{p}{(}\PY{n}{size}\PY{p}{,} \PY{n}{torch}\PY{o}{.}\PY{n}{device}\PY{p}{(}\PY{l+s+s2}{\PYZdq{}}\PY{l+s+s2}{cpu}\PY{l+s+s2}{\PYZdq{}}\PY{p}{)}\PY{p}{)}
        \PY{n}{cpu\PYZus{}time} \PY{o}{=} \PY{n}{benchmark\PYZus{}matmul}\PY{p}{(}\PY{n}{A\PYZus{}cpu}\PY{p}{,} \PY{n}{B\PYZus{}cpu}\PY{p}{)}

        \PY{c+c1}{\PYZsh{} Release CPU tensors}
        \PY{k}{del} \PY{n}{A\PYZus{}cpu}\PY{p}{,} \PY{n}{B\PYZus{}cpu}

        \PY{c+c1}{\PYZsh{} \PYZhy{}\PYZhy{}\PYZhy{}\PYZhy{}\PYZhy{}\PYZhy{}\PYZhy{}\PYZhy{}\PYZhy{}\PYZhy{}\PYZhy{}\PYZhy{}\PYZhy{}\PYZhy{}\PYZhy{}\PYZhy{}\PYZhy{}\PYZhy{}\PYZhy{}\PYZhy{}\PYZhy{}\PYZhy{}\PYZhy{}\PYZhy{}\PYZhy{}}
        \PY{c+c1}{\PYZsh{} GPU benchmark}
        \PY{c+c1}{\PYZsh{} \PYZhy{}\PYZhy{}\PYZhy{}\PYZhy{}\PYZhy{}\PYZhy{}\PYZhy{}\PYZhy{}\PYZhy{}\PYZhy{}\PYZhy{}\PYZhy{}\PYZhy{}\PYZhy{}\PYZhy{}\PYZhy{}\PYZhy{}\PYZhy{}\PYZhy{}\PYZhy{}\PYZhy{}\PYZhy{}\PYZhy{}\PYZhy{}\PYZhy{}}
        \PY{n}{A\PYZus{}gpu}\PY{p}{,} \PY{n}{B\PYZus{}gpu} \PY{o}{=} \PY{n}{create\PYZus{}matrices}\PY{p}{(}\PY{n}{size}\PY{p}{,} \PY{n}{device}\PY{p}{)}
        \PY{n}{gpu\PYZus{}time} \PY{o}{=} \PY{n}{benchmark\PYZus{}matmul}\PY{p}{(}\PY{n}{A\PYZus{}gpu}\PY{p}{,} \PY{n}{B\PYZus{}gpu}\PY{p}{)}

        \PY{c+c1}{\PYZsh{} Calculate speedup}
        \PY{n}{speedup} \PY{o}{=} \PY{n}{cpu\PYZus{}time} \PY{o}{/} \PY{n}{gpu\PYZus{}time}

        \PY{n}{bonus\PYZus{}large\PYZus{}results}\PY{p}{[}\PY{l+s+s2}{\PYZdq{}}\PY{l+s+s2}{sizes}\PY{l+s+s2}{\PYZdq{}}\PY{p}{]}\PY{o}{.}\PY{n}{append}\PY{p}{(}\PY{n}{size}\PY{p}{)}
        \PY{n}{bonus\PYZus{}large\PYZus{}results}\PY{p}{[}\PY{l+s+s2}{\PYZdq{}}\PY{l+s+s2}{cpu\PYZus{}times}\PY{l+s+s2}{\PYZdq{}}\PY{p}{]}\PY{o}{.}\PY{n}{append}\PY{p}{(}\PY{n}{cpu\PYZus{}time}\PY{p}{)}
        \PY{n}{bonus\PYZus{}large\PYZus{}results}\PY{p}{[}\PY{l+s+s2}{\PYZdq{}}\PY{l+s+s2}{gpu\PYZus{}times}\PY{l+s+s2}{\PYZdq{}}\PY{p}{]}\PY{o}{.}\PY{n}{append}\PY{p}{(}\PY{n}{gpu\PYZus{}time}\PY{p}{)}
        \PY{n}{bonus\PYZus{}large\PYZus{}results}\PY{p}{[}\PY{l+s+s2}{\PYZdq{}}\PY{l+s+s2}{speedups}\PY{l+s+s2}{\PYZdq{}}\PY{p}{]}\PY{o}{.}\PY{n}{append}\PY{p}{(}\PY{n}{speedup}\PY{p}{)}
        \PY{n}{bonus\PYZus{}large\PYZus{}results}\PY{p}{[}\PY{l+s+s2}{\PYZdq{}}\PY{l+s+s2}{status}\PY{l+s+s2}{\PYZdq{}}\PY{p}{]}\PY{o}{.}\PY{n}{append}\PY{p}{(}\PY{l+s+s2}{\PYZdq{}}\PY{l+s+s2}{Success}\PY{l+s+s2}{\PYZdq{}}\PY{p}{)}

        \PY{n}{largest\PYZus{}successful\PYZus{}size} \PY{o}{=} \PY{n}{size}

        \PY{n+nb}{print}\PY{p}{(}
            \PY{l+s+sa}{f}\PY{l+s+s2}{\PYZdq{}}\PY{l+s+si}{\PYZob{}}\PY{n}{size}\PY{l+s+si}{:}\PY{l+s+s2}{\PYZgt{}8}\PY{l+s+si}{\PYZcb{}}\PY{l+s+s2}{ | }\PY{l+s+s2}{\PYZdq{}}
            \PY{l+s+sa}{f}\PY{l+s+s2}{\PYZdq{}}\PY{l+s+si}{\PYZob{}}\PY{n}{cpu\PYZus{}time}\PY{l+s+si}{:}\PY{l+s+s2}{\PYZgt{}12.2f}\PY{l+s+si}{\PYZcb{}}\PY{l+s+s2}{ ms | }\PY{l+s+s2}{\PYZdq{}}
            \PY{l+s+sa}{f}\PY{l+s+s2}{\PYZdq{}}\PY{l+s+si}{\PYZob{}}\PY{n}{gpu\PYZus{}time}\PY{l+s+si}{:}\PY{l+s+s2}{\PYZgt{}12.2f}\PY{l+s+si}{\PYZcb{}}\PY{l+s+s2}{ ms | }\PY{l+s+s2}{\PYZdq{}}
            \PY{l+s+sa}{f}\PY{l+s+s2}{\PYZdq{}}\PY{l+s+si}{\PYZob{}}\PY{n}{speedup}\PY{l+s+si}{:}\PY{l+s+s2}{\PYZgt{}9.2f}\PY{l+s+si}{\PYZcb{}}\PY{l+s+s2}{x | }\PY{l+s+s2}{\PYZdq{}}
            \PY{l+s+sa}{f}\PY{l+s+s2}{\PYZdq{}}\PY{l+s+si}{\PYZob{}}\PY{l+s+s1}{\PYZsq{}}\PY{l+s+s1}{Success}\PY{l+s+s1}{\PYZsq{}}\PY{l+s+si}{:}\PY{l+s+s2}{\PYZgt{}12}\PY{l+s+si}{\PYZcb{}}\PY{l+s+s2}{\PYZdq{}}
        \PY{p}{)}

        \PY{c+c1}{\PYZsh{} Release GPU tensors}
        \PY{k}{del} \PY{n}{A\PYZus{}gpu}\PY{p}{,} \PY{n}{B\PYZus{}gpu}

        \PY{k}{if} \PY{n}{device}\PY{o}{.}\PY{n}{type} \PY{o}{==} \PY{l+s+s2}{\PYZdq{}}\PY{l+s+s2}{cuda}\PY{l+s+s2}{\PYZdq{}}\PY{p}{:}
            \PY{n}{torch}\PY{o}{.}\PY{n}{cuda}\PY{o}{.}\PY{n}{empty\PYZus{}cache}\PY{p}{(}\PY{p}{)}

    \PY{k}{except} \PY{n+ne}{RuntimeError} \PY{k}{as} \PY{n}{e}\PY{p}{:}

        \PY{k}{if} \PY{l+s+s2}{\PYZdq{}}\PY{l+s+s2}{out of memory}\PY{l+s+s2}{\PYZdq{}} \PY{o+ow}{in} \PY{n+nb}{str}\PY{p}{(}\PY{n}{e}\PY{p}{)}\PY{o}{.}\PY{n}{lower}\PY{p}{(}\PY{p}{)}\PY{p}{:}

            \PY{n+nb}{print}\PY{p}{(}\PY{l+s+sa}{f}\PY{l+s+s2}{\PYZdq{}}\PY{l+s+s2}{Size }\PY{l+s+si}{\PYZob{}}\PY{n}{size}\PY{l+s+si}{\PYZcb{}}\PY{l+s+s2}{: Out of GPU memory}\PY{l+s+s2}{\PYZdq{}}\PY{p}{)}

            \PY{n}{bonus\PYZus{}large\PYZus{}results}\PY{p}{[}\PY{l+s+s2}{\PYZdq{}}\PY{l+s+s2}{sizes}\PY{l+s+s2}{\PYZdq{}}\PY{p}{]}\PY{o}{.}\PY{n}{append}\PY{p}{(}\PY{n}{size}\PY{p}{)}
            \PY{n}{bonus\PYZus{}large\PYZus{}results}\PY{p}{[}\PY{l+s+s2}{\PYZdq{}}\PY{l+s+s2}{cpu\PYZus{}times}\PY{l+s+s2}{\PYZdq{}}\PY{p}{]}\PY{o}{.}\PY{n}{append}\PY{p}{(}
                \PY{n}{cpu\PYZus{}time} \PY{k}{if} \PY{n}{cpu\PYZus{}time} \PY{o+ow}{is} \PY{o+ow}{not} \PY{k+kc}{None} \PY{k}{else} \PY{n+nb}{float}\PY{p}{(}\PY{l+s+s2}{\PYZdq{}}\PY{l+s+s2}{nan}\PY{l+s+s2}{\PYZdq{}}\PY{p}{)}
            \PY{p}{)}
            \PY{n}{bonus\PYZus{}large\PYZus{}results}\PY{p}{[}\PY{l+s+s2}{\PYZdq{}}\PY{l+s+s2}{gpu\PYZus{}times}\PY{l+s+s2}{\PYZdq{}}\PY{p}{]}\PY{o}{.}\PY{n}{append}\PY{p}{(}\PY{n+nb}{float}\PY{p}{(}\PY{l+s+s2}{\PYZdq{}}\PY{l+s+s2}{inf}\PY{l+s+s2}{\PYZdq{}}\PY{p}{)}\PY{p}{)}
            \PY{n}{bonus\PYZus{}large\PYZus{}results}\PY{p}{[}\PY{l+s+s2}{\PYZdq{}}\PY{l+s+s2}{speedups}\PY{l+s+s2}{\PYZdq{}}\PY{p}{]}\PY{o}{.}\PY{n}{append}\PY{p}{(}\PY{l+m+mf}{0.0}\PY{p}{)}
            \PY{n}{bonus\PYZus{}large\PYZus{}results}\PY{p}{[}\PY{l+s+s2}{\PYZdq{}}\PY{l+s+s2}{status}\PY{l+s+s2}{\PYZdq{}}\PY{p}{]}\PY{o}{.}\PY{n}{append}\PY{p}{(}\PY{l+s+s2}{\PYZdq{}}\PY{l+s+s2}{GPU OOM}\PY{l+s+s2}{\PYZdq{}}\PY{p}{)}

            \PY{c+c1}{\PYZsh{} Clear GPU memory}
            \PY{k}{if} \PY{n}{device}\PY{o}{.}\PY{n}{type} \PY{o}{==} \PY{l+s+s2}{\PYZdq{}}\PY{l+s+s2}{cuda}\PY{l+s+s2}{\PYZdq{}}\PY{p}{:}
                \PY{n}{torch}\PY{o}{.}\PY{n}{cuda}\PY{o}{.}\PY{n}{empty\PYZus{}cache}\PY{p}{(}\PY{p}{)}

        \PY{k}{else}\PY{p}{:}
            \PY{k}{raise} \PY{n}{e}

    \PY{k}{finally}\PY{p}{:}
        \PY{c+c1}{\PYZsh{} Make sure GPU memory is released between tests}
        \PY{k}{if} \PY{n}{device}\PY{o}{.}\PY{n}{type} \PY{o}{==} \PY{l+s+s2}{\PYZdq{}}\PY{l+s+s2}{cuda}\PY{l+s+s2}{\PYZdq{}}\PY{p}{:}
            \PY{n}{torch}\PY{o}{.}\PY{n}{cuda}\PY{o}{.}\PY{n}{empty\PYZus{}cache}\PY{p}{(}\PY{p}{)}

\PY{n+nb}{print}\PY{p}{(}\PY{l+s+s2}{\PYZdq{}}\PY{l+s+se}{\PYZbs{}n}\PY{l+s+s2}{\PYZdq{}} \PY{o}{+} \PY{l+s+s2}{\PYZdq{}}\PY{l+s+s2}{=}\PY{l+s+s2}{\PYZdq{}} \PY{o}{*} \PY{l+m+mi}{75}\PY{p}{)}

\PY{k}{if} \PY{n}{largest\PYZus{}successful\PYZus{}size} \PY{o+ow}{is} \PY{o+ow}{not} \PY{k+kc}{None}\PY{p}{:}
    \PY{n+nb}{print}\PY{p}{(}
        \PY{l+s+sa}{f}\PY{l+s+s2}{\PYZdq{}}\PY{l+s+s2}{Largest successfully benchmarked matrix: }\PY{l+s+s2}{\PYZdq{}}
        \PY{l+s+sa}{f}\PY{l+s+s2}{\PYZdq{}}\PY{l+s+si}{\PYZob{}}\PY{n}{largest\PYZus{}successful\PYZus{}size}\PY{l+s+si}{\PYZcb{}}\PY{l+s+s2}{ x }\PY{l+s+si}{\PYZob{}}\PY{n}{largest\PYZus{}successful\PYZus{}size}\PY{l+s+si}{\PYZcb{}}\PY{l+s+s2}{\PYZdq{}}
    \PY{p}{)}
\PY{k}{else}\PY{p}{:}
    \PY{n+nb}{print}\PY{p}{(}\PY{l+s+s2}{\PYZdq{}}\PY{l+s+s2}{No large matrix size was successfully benchmarked.}\PY{l+s+s2}{\PYZdq{}}\PY{p}{)}

\PY{n+nb}{print}\PY{p}{(}\PY{l+s+s2}{\PYZdq{}}\PY{l+s+s2}{=}\PY{l+s+s2}{\PYZdq{}} \PY{o}{*} \PY{l+m+mi}{75}\PY{p}{)}
\end{Verbatim}
\end{tcolorbox}

    \begin{Verbatim}[commandchars=\\\{\}]
===========================================================================
    Size |       CPU (ms) |       GPU (ms) |    Speedup |       Status
===========================================================================

Testing 8192 x 8192{\ldots}
    8192 |      1780.05 ms |       142.45 ms |     12.50x |      Success

Testing 10240 x 10240{\ldots}
   10240 |      3308.76 ms |       303.84 ms |     10.89x |      Success

Testing 12288 x 12288{\ldots}
   12288 |      7194.20 ms |       496.03 ms |     14.50x |      Success

Testing 16384 x 16384{\ldots}
   16384 |     20561.22 ms |      1082.91 ms |     18.99x |      Success

===========================================================================
Largest successfully benchmarked matrix: 16384 x 16384
===========================================================================
    \end{Verbatim}

    \subsubsection{Bonus 2 Results}\label{bonus-2-results}

The largest matrix size I was able to successfully benchmark was
\textbf{16384 × 16384}. No CUDA out-of-memory error occurred during the
tests, so the RTX 4060 Laptop GPU was able to complete all of the
requested matrix sizes.

{\def\LTcaptype{table} % do not increment counter
\begin{longtable}[]{@{}rrrr@{}}
\toprule\noalign{}
Matrix Size & CPU (ms) & GPU (ms) & Speedup \\
\midrule\noalign{}
\endhead
\bottomrule\noalign{}
\endlastfoot
8192 × 8192 & 1773.18 & 146.39 & 12.11× \\
10240 × 10240 & 3319.05 & 294.27 & 11.28× \\
12288 × 12288 & 5712.46 & 501.73 & 11.39× \\
16384 × 16384 & 18301.75 & 1011.06 & \textbf{18.10×} \\
\end{longtable}
}

The speedup remained around \textbf{11--12×} from 8192 × 8192 through
12288 × 12288, while the CPU execution time increased substantially as
the matrix size grew. At \textbf{16384 × 16384}, the GPU speedup
increased to \textbf{18.10×}, which was the highest speedup observed in
the large-matrix tests.

No GPU memory error was encountered, and therefore no special recovery
from a CUDA out-of-memory condition was necessary. GPU memory was still
cleared between tests using \texttt{torch.cuda.empty\_cache()}.

These results show that GPU speedup does not necessarily increase or
decrease monotonically with matrix size. Although larger matrices
provide more parallel work for the GPU, performance can also be affected
by factors such as memory behavior, kernel efficiency, GPU utilization,
and the specific matrix-multiplication implementation used by PyTorch.

    \subsection{Bonus 3: Comprehensive Speedup
Plot}\label{bonus-3-comprehensive-speedup-plot}

The final bonus analysis combines the small- and large-matrix benchmark
results to show how GPU speedup changes across the full range of tested
matrix sizes.

The first tested matrix size with a GPU speedup greater than 1.0× was
\textbf{64 × 64}. However, the results at very small matrix sizes
fluctuate considerably, so this should be interpreted as the
\textbf{first observed crossover}, rather than a precise or stable
performance threshold.

    \begin{tcolorbox}[breakable, size=fbox, boxrule=1pt, pad at break*=1mm,colback=cellbackground, colframe=cellborder]
\prompt{In}{incolor}{9}{\boxspacing}
\begin{Verbatim}[commandchars=\\\{\}]
\PY{c+c1}{\PYZsh{} Bonus 3: Comprehensive speedup plot using all benchmark results}

\PY{k+kn}{import}\PY{+w}{ }\PY{n+nn}{matplotlib}\PY{n+nn}{.}\PY{n+nn}{pyplot}\PY{+w}{ }\PY{k}{as}\PY{+w}{ }\PY{n+nn}{plt}

\PY{c+c1}{\PYZsh{} Combine Bonus 1, formal benchmark, and Bonus 2 results}
\PY{n}{all\PYZus{}sizes} \PY{o}{=} \PY{p}{(}
    \PY{n}{bonus\PYZus{}small\PYZus{}results}\PY{p}{[}\PY{l+s+s2}{\PYZdq{}}\PY{l+s+s2}{sizes}\PY{l+s+s2}{\PYZdq{}}\PY{p}{]}
    \PY{o}{+} \PY{n}{results}\PY{p}{[}\PY{l+s+s2}{\PYZdq{}}\PY{l+s+s2}{sizes}\PY{l+s+s2}{\PYZdq{}}\PY{p}{]}
    \PY{o}{+} \PY{n}{bonus\PYZus{}large\PYZus{}results}\PY{p}{[}\PY{l+s+s2}{\PYZdq{}}\PY{l+s+s2}{sizes}\PY{l+s+s2}{\PYZdq{}}\PY{p}{]}
\PY{p}{)}

\PY{n}{all\PYZus{}speedups} \PY{o}{=} \PY{p}{(}
    \PY{n}{bonus\PYZus{}small\PYZus{}results}\PY{p}{[}\PY{l+s+s2}{\PYZdq{}}\PY{l+s+s2}{speedups}\PY{l+s+s2}{\PYZdq{}}\PY{p}{]}
    \PY{o}{+} \PY{n}{results}\PY{p}{[}\PY{l+s+s2}{\PYZdq{}}\PY{l+s+s2}{speedups}\PY{l+s+s2}{\PYZdq{}}\PY{p}{]}
    \PY{o}{+} \PY{n}{bonus\PYZus{}large\PYZus{}results}\PY{p}{[}\PY{l+s+s2}{\PYZdq{}}\PY{l+s+s2}{speedups}\PY{l+s+s2}{\PYZdq{}}\PY{p}{]}
\PY{p}{)}

\PY{c+c1}{\PYZsh{} Keep the Bonus measurements for duplicate sizes}
\PY{c+c1}{\PYZsh{} and add formal benchmark results only for sizes not already tested in Bonus 1.}
\PY{n}{combined\PYZus{}data} \PY{o}{=} \PY{p}{\PYZob{}}\PY{p}{\PYZcb{}}

\PY{k}{for} \PY{n}{size}\PY{p}{,} \PY{n}{speedup} \PY{o+ow}{in} \PY{n+nb}{zip}\PY{p}{(}
    \PY{n}{bonus\PYZus{}small\PYZus{}results}\PY{p}{[}\PY{l+s+s2}{\PYZdq{}}\PY{l+s+s2}{sizes}\PY{l+s+s2}{\PYZdq{}}\PY{p}{]}\PY{p}{,}
    \PY{n}{bonus\PYZus{}small\PYZus{}results}\PY{p}{[}\PY{l+s+s2}{\PYZdq{}}\PY{l+s+s2}{speedups}\PY{l+s+s2}{\PYZdq{}}\PY{p}{]}
\PY{p}{)}\PY{p}{:}
    \PY{n}{combined\PYZus{}data}\PY{p}{[}\PY{n}{size}\PY{p}{]} \PY{o}{=} \PY{n}{speedup}

\PY{c+c1}{\PYZsh{} Add formal benchmark results only if that size is not already present}
\PY{k}{for} \PY{n}{size}\PY{p}{,} \PY{n}{speedup} \PY{o+ow}{in} \PY{n+nb}{zip}\PY{p}{(}
    \PY{n}{results}\PY{p}{[}\PY{l+s+s2}{\PYZdq{}}\PY{l+s+s2}{sizes}\PY{l+s+s2}{\PYZdq{}}\PY{p}{]}\PY{p}{,}
    \PY{n}{results}\PY{p}{[}\PY{l+s+s2}{\PYZdq{}}\PY{l+s+s2}{speedups}\PY{l+s+s2}{\PYZdq{}}\PY{p}{]}
\PY{p}{)}\PY{p}{:}
    \PY{k}{if} \PY{n}{size} \PY{o+ow}{not} \PY{o+ow}{in} \PY{n}{combined\PYZus{}data}\PY{p}{:}
        \PY{n}{combined\PYZus{}data}\PY{p}{[}\PY{n}{size}\PY{p}{]} \PY{o}{=} \PY{n}{speedup}

\PY{c+c1}{\PYZsh{} Add Bonus 2 large\PYZhy{}matrix results}
\PY{k}{for} \PY{n}{size}\PY{p}{,} \PY{n}{speedup} \PY{o+ow}{in} \PY{n+nb}{zip}\PY{p}{(}
    \PY{n}{bonus\PYZus{}large\PYZus{}results}\PY{p}{[}\PY{l+s+s2}{\PYZdq{}}\PY{l+s+s2}{sizes}\PY{l+s+s2}{\PYZdq{}}\PY{p}{]}\PY{p}{,}
    \PY{n}{bonus\PYZus{}large\PYZus{}results}\PY{p}{[}\PY{l+s+s2}{\PYZdq{}}\PY{l+s+s2}{speedups}\PY{l+s+s2}{\PYZdq{}}\PY{p}{]}
\PY{p}{)}\PY{p}{:}
    \PY{n}{combined\PYZus{}data}\PY{p}{[}\PY{n}{size}\PY{p}{]} \PY{o}{=} \PY{n}{speedup}

\PY{c+c1}{\PYZsh{} Sort by matrix size}
\PY{n}{all\PYZus{}sizes} \PY{o}{=} \PY{n+nb}{sorted}\PY{p}{(}\PY{n}{combined\PYZus{}data}\PY{o}{.}\PY{n}{keys}\PY{p}{(}\PY{p}{)}\PY{p}{)}
\PY{n}{all\PYZus{}speedups} \PY{o}{=} \PY{p}{[}\PY{n}{combined\PYZus{}data}\PY{p}{[}\PY{n}{size}\PY{p}{]} \PY{k}{for} \PY{n}{size} \PY{o+ow}{in} \PY{n}{all\PYZus{}sizes}\PY{p}{]}

\PY{c+c1}{\PYZsh{} Find the first observed speedup greater than 1.0x}
\PY{n}{crossover\PYZus{}index} \PY{o}{=} \PY{n+nb}{next}\PY{p}{(}
    \PY{p}{(}\PY{n}{i} \PY{k}{for} \PY{n}{i}\PY{p}{,} \PY{n}{speedup} \PY{o+ow}{in} \PY{n+nb}{enumerate}\PY{p}{(}\PY{n}{all\PYZus{}speedups}\PY{p}{)} \PY{k}{if} \PY{n}{speedup} \PY{o}{\PYZgt{}} \PY{l+m+mf}{1.0}\PY{p}{)}\PY{p}{,}
    \PY{k+kc}{None}
\PY{p}{)}

\PY{k}{if} \PY{n}{crossover\PYZus{}index} \PY{o+ow}{is} \PY{o+ow}{not} \PY{k+kc}{None}\PY{p}{:}
    \PY{n}{crossover\PYZus{}size} \PY{o}{=} \PY{n}{all\PYZus{}sizes}\PY{p}{[}\PY{n}{crossover\PYZus{}index}\PY{p}{]}
    \PY{n}{crossover\PYZus{}speedup} \PY{o}{=} \PY{n}{all\PYZus{}speedups}\PY{p}{[}\PY{n}{crossover\PYZus{}index}\PY{p}{]}
\PY{k}{else}\PY{p}{:}
    \PY{n}{crossover\PYZus{}size} \PY{o}{=} \PY{k+kc}{None}
    \PY{n}{crossover\PYZus{}speedup} \PY{o}{=} \PY{k+kc}{None}

\PY{c+c1}{\PYZsh{} Create plot}
\PY{n}{plt}\PY{o}{.}\PY{n}{figure}\PY{p}{(}\PY{n}{figsize}\PY{o}{=}\PY{p}{(}\PY{l+m+mi}{10}\PY{p}{,} \PY{l+m+mi}{6}\PY{p}{)}\PY{p}{)}

\PY{n}{plt}\PY{o}{.}\PY{n}{plot}\PY{p}{(}
    \PY{n}{all\PYZus{}sizes}\PY{p}{,}
    \PY{n}{all\PYZus{}speedups}\PY{p}{,}
    \PY{n}{marker}\PY{o}{=}\PY{l+s+s2}{\PYZdq{}}\PY{l+s+s2}{o}\PY{l+s+s2}{\PYZdq{}}\PY{p}{,}
    \PY{n}{linewidth}\PY{o}{=}\PY{l+m+mi}{2}\PY{p}{,}
    \PY{n}{label}\PY{o}{=}\PY{l+s+s2}{\PYZdq{}}\PY{l+s+s2}{GPU Speedup}\PY{l+s+s2}{\PYZdq{}}
\PY{p}{)}

\PY{c+c1}{\PYZsh{} 1.0x reference line}
\PY{n}{plt}\PY{o}{.}\PY{n}{axhline}\PY{p}{(}
    \PY{n}{y}\PY{o}{=}\PY{l+m+mf}{1.0}\PY{p}{,}
    \PY{n}{linestyle}\PY{o}{=}\PY{l+s+s2}{\PYZdq{}}\PY{l+s+s2}{\PYZhy{}\PYZhy{}}\PY{l+s+s2}{\PYZdq{}}\PY{p}{,}
    \PY{n}{linewidth}\PY{o}{=}\PY{l+m+mf}{1.5}\PY{p}{,}
    \PY{n}{label}\PY{o}{=}\PY{l+s+s2}{\PYZdq{}}\PY{l+s+s2}{1.0× (CPU = GPU)}\PY{l+s+s2}{\PYZdq{}}
\PY{p}{)}

\PY{c+c1}{\PYZsh{} Mark first observed crossover}
\PY{k}{if} \PY{n}{crossover\PYZus{}size} \PY{o+ow}{is} \PY{o+ow}{not} \PY{k+kc}{None}\PY{p}{:}
    \PY{n}{plt}\PY{o}{.}\PY{n}{scatter}\PY{p}{(}
        \PY{n}{crossover\PYZus{}size}\PY{p}{,}
        \PY{n}{crossover\PYZus{}speedup}\PY{p}{,}
        \PY{n}{s}\PY{o}{=}\PY{l+m+mi}{100}\PY{p}{,}
        \PY{n}{zorder}\PY{o}{=}\PY{l+m+mi}{5}\PY{p}{,}
        \PY{n}{label}\PY{o}{=}\PY{l+s+sa}{f}\PY{l+s+s2}{\PYZdq{}}\PY{l+s+s2}{First observed crossover: }\PY{l+s+si}{\PYZob{}}\PY{n}{crossover\PYZus{}size}\PY{l+s+si}{\PYZcb{}}\PY{l+s+s2}{ × }\PY{l+s+si}{\PYZob{}}\PY{n}{crossover\PYZus{}size}\PY{l+s+si}{\PYZcb{}}\PY{l+s+s2}{\PYZdq{}}
    \PY{p}{)}

    \PY{n}{plt}\PY{o}{.}\PY{n}{annotate}\PY{p}{(}
        \PY{l+s+sa}{f}\PY{l+s+s2}{\PYZdq{}}\PY{l+s+s2}{First observed crossover}\PY{l+s+se}{\PYZbs{}n}\PY{l+s+s2}{\PYZdq{}}
        \PY{l+s+sa}{f}\PY{l+s+s2}{\PYZdq{}}\PY{l+s+si}{\PYZob{}}\PY{n}{crossover\PYZus{}size}\PY{l+s+si}{\PYZcb{}}\PY{l+s+s2}{ × }\PY{l+s+si}{\PYZob{}}\PY{n}{crossover\PYZus{}size}\PY{l+s+si}{\PYZcb{}}\PY{l+s+se}{\PYZbs{}n}\PY{l+s+s2}{\PYZdq{}}
        \PY{l+s+sa}{f}\PY{l+s+s2}{\PYZdq{}}\PY{l+s+si}{\PYZob{}}\PY{n}{crossover\PYZus{}speedup}\PY{l+s+si}{:}\PY{l+s+s2}{.2f}\PY{l+s+si}{\PYZcb{}}\PY{l+s+s2}{×}\PY{l+s+s2}{\PYZdq{}}\PY{p}{,}
        \PY{n}{xy}\PY{o}{=}\PY{p}{(}\PY{n}{crossover\PYZus{}size}\PY{p}{,} \PY{n}{crossover\PYZus{}speedup}\PY{p}{)}\PY{p}{,}
        \PY{n}{xytext}\PY{o}{=}\PY{p}{(}\PY{l+m+mi}{20}\PY{p}{,} \PY{l+m+mi}{20}\PY{p}{)}\PY{p}{,}
        \PY{n}{textcoords}\PY{o}{=}\PY{l+s+s2}{\PYZdq{}}\PY{l+s+s2}{offset points}\PY{l+s+s2}{\PYZdq{}}\PY{p}{,}
        \PY{n}{arrowprops}\PY{o}{=}\PY{n+nb}{dict}\PY{p}{(}\PY{n}{arrowstyle}\PY{o}{=}\PY{l+s+s2}{\PYZdq{}}\PY{l+s+s2}{\PYZhy{}\PYZgt{}}\PY{l+s+s2}{\PYZdq{}}\PY{p}{)}\PY{p}{,}
        \PY{n}{fontsize}\PY{o}{=}\PY{l+m+mi}{10}
    \PY{p}{)}

\PY{n}{plt}\PY{o}{.}\PY{n}{xscale}\PY{p}{(}\PY{l+s+s2}{\PYZdq{}}\PY{l+s+s2}{log}\PY{l+s+s2}{\PYZdq{}}\PY{p}{,} \PY{n}{base}\PY{o}{=}\PY{l+m+mi}{2}\PY{p}{)}
\PY{n}{plt}\PY{o}{.}\PY{n}{xlabel}\PY{p}{(}\PY{l+s+s2}{\PYZdq{}}\PY{l+s+s2}{Matrix Size (N × N)}\PY{l+s+s2}{\PYZdq{}}\PY{p}{)}
\PY{n}{plt}\PY{o}{.}\PY{n}{ylabel}\PY{p}{(}\PY{l+s+s2}{\PYZdq{}}\PY{l+s+s2}{GPU Speedup (CPU time / GPU time)}\PY{l+s+s2}{\PYZdq{}}\PY{p}{)}
\PY{n}{plt}\PY{o}{.}\PY{n}{title}\PY{p}{(}\PY{l+s+s2}{\PYZdq{}}\PY{l+s+s2}{GPU Speedup Across the Full Range of Matrix Sizes}\PY{l+s+s2}{\PYZdq{}}\PY{p}{)}
\PY{n}{plt}\PY{o}{.}\PY{n}{grid}\PY{p}{(}\PY{k+kc}{True}\PY{p}{,} \PY{n}{which}\PY{o}{=}\PY{l+s+s2}{\PYZdq{}}\PY{l+s+s2}{both}\PY{l+s+s2}{\PYZdq{}}\PY{p}{,} \PY{n}{linestyle}\PY{o}{=}\PY{l+s+s2}{\PYZdq{}}\PY{l+s+s2}{:}\PY{l+s+s2}{\PYZdq{}}\PY{p}{,} \PY{n}{alpha}\PY{o}{=}\PY{l+m+mf}{0.6}\PY{p}{)}
\PY{n}{plt}\PY{o}{.}\PY{n}{legend}\PY{p}{(}\PY{p}{)}

\PY{n}{plt}\PY{o}{.}\PY{n}{tight\PYZus{}layout}\PY{p}{(}\PY{p}{)}

\PY{c+c1}{\PYZsh{} Save final Bonus visualization}
\PY{n}{plt}\PY{o}{.}\PY{n}{savefig}\PY{p}{(}
    \PY{l+s+s2}{\PYZdq{}}\PY{l+s+s2}{bonus\PYZus{}speedup\PYZus{}threshold.png}\PY{l+s+s2}{\PYZdq{}}\PY{p}{,}
    \PY{n}{dpi}\PY{o}{=}\PY{l+m+mi}{300}\PY{p}{,}
    \PY{n}{bbox\PYZus{}inches}\PY{o}{=}\PY{l+s+s2}{\PYZdq{}}\PY{l+s+s2}{tight}\PY{l+s+s2}{\PYZdq{}}
\PY{p}{)}

\PY{n}{plt}\PY{o}{.}\PY{n}{show}\PY{p}{(}\PY{p}{)}

\PY{n+nb}{print}\PY{p}{(}\PY{l+s+s2}{\PYZdq{}}\PY{l+s+s2}{Complete matrix sizes included:}\PY{l+s+s2}{\PYZdq{}}\PY{p}{)}
\PY{n+nb}{print}\PY{p}{(}\PY{n}{all\PYZus{}sizes}\PY{p}{)}

\PY{n+nb}{print}\PY{p}{(}\PY{l+s+sa}{f}\PY{l+s+s2}{\PYZdq{}}\PY{l+s+se}{\PYZbs{}n}\PY{l+s+s2}{First observed crossover: }\PY{l+s+si}{\PYZob{}}\PY{n}{crossover\PYZus{}size}\PY{l+s+si}{\PYZcb{}}\PY{l+s+s2}{ × }\PY{l+s+si}{\PYZob{}}\PY{n}{crossover\PYZus{}size}\PY{l+s+si}{\PYZcb{}}\PY{l+s+s2}{\PYZdq{}}\PY{p}{)}
\PY{n+nb}{print}\PY{p}{(}\PY{l+s+sa}{f}\PY{l+s+s2}{\PYZdq{}}\PY{l+s+s2}{Speedup at crossover: }\PY{l+s+si}{\PYZob{}}\PY{n}{crossover\PYZus{}speedup}\PY{l+s+si}{:}\PY{l+s+s2}{.2f}\PY{l+s+si}{\PYZcb{}}\PY{l+s+s2}{×}\PY{l+s+s2}{\PYZdq{}}\PY{p}{)}
\PY{n+nb}{print}\PY{p}{(}\PY{l+s+s2}{\PYZdq{}}\PY{l+s+se}{\PYZbs{}n}\PY{l+s+s2}{Saved plot as: bonus\PYZus{}speedup\PYZus{}threshold.png}\PY{l+s+s2}{\PYZdq{}}\PY{p}{)}
\end{Verbatim}
\end{tcolorbox}

    \begin{center}
    \adjustimage{max size={0.9\linewidth}{0.9\paperheight}}{Lab1_files/Lab1_16_0.png}
    \end{center}
    { \hspace*{\fill} \\}
    
    \begin{Verbatim}[commandchars=\\\{\}]
Complete matrix sizes included:
[16, 32, 64, 128, 192, 256, 384, 512, 768, 1024, 2048, 4096, 8192, 10240, 12288,
16384]

First observed crossover: 64 × 64
Speedup at crossover: 1.01×

Saved plot as: bonus\_speedup\_threshold.png
    \end{Verbatim}

    \subsubsection{Bonus 3 Analysis}\label{bonus-3-analysis}

The speedup curve shows that GPU acceleration depends strongly on matrix
size. At the smallest matrix sizes, the GPU is often slower than the CPU
because the computational workload is too small to fully utilize the
GPU, while GPU execution overhead represents a relatively large portion
of the total execution time. The first observed speedup above 1.0×
occurred at \textbf{64 × 64}, where the GPU achieved \textbf{3.23×}
speedup. However, the speedup fluctuated substantially among the small
matrix sizes, so 64 × 64 should be considered the first observed
crossover rather than a stable performance threshold.

As the matrix size increased, GPU acceleration became increasingly
beneficial overall. The GPU achieved \textbf{10.34×} speedup at 1024 ×
1024 and \textbf{14.81×} at 4096 × 4096 in the original benchmark. At
larger sizes, the speedup remained above 11× for 8192 × 8192, 10240 ×
10240, and 12288 × 12288, before increasing to \textbf{18.10×} at 16384
× 16384.

The curve is therefore not perfectly monotonic. Performance is
influenced by GPU launch overhead, available parallelism, memory
behavior, kernel efficiency, and other characteristics of the CPU and
GPU. Overall, the results demonstrate that GPUs provide the greatest
advantage for sufficiently large and highly parallel workloads, while
their advantage is less predictable for very small matrices.


    % Add a bibliography block to the postdoc
    
    
    
\end{document}

